% Chapter 12. All quoted numbers come from chapters/ch12-numbers.tex, written by code/ch12/run_all.py.
% Generated by code/ch12/run_all.py. Do not edit by hand.
\newcommand{\bvLassoBias}{0.209}
\newcommand{\bvLassoCut}{79}
\newcommand{\bvLassoErr}{0.576}
\newcommand{\bvLassoLam}{0.200}
\newcommand{\bvLassoNz}{11.9}
\newcommand{\bvLassoVar}{0.368}
\newcommand{\bvN}{100}
\newcommand{\bvNonzero}{5}
\newcommand{\bvNte}{4,000}
\newcommand{\bvOlsBias}{0.011}
\newcommand{\bvOlsErr}{2.800}
\newcommand{\bvOlsVar}{2.789}
\newcommand{\bvP}{40}
\newcommand{\bvPathNzOne}{10}
\newcommand{\bvReps}{300}
\newcommand{\bvRho}{0.6}
\newcommand{\bvRidgeBias}{0.416}
\newcommand{\bvRidgeCut}{68}
\newcommand{\bvRidgeDf}{16.8}
\newcommand{\bvRidgeErr}{0.892}
\newcommand{\bvRidgeLam}{75.0}
\newcommand{\bvRidgeSmall}{2.766}
\newcommand{\bvRidgeSmallLam}{0.1}
\newcommand{\bvRidgeVar}{0.476}
\newcommand{\bvSig}{2}
\newcommand{\bvSnr}{0.54}
\newcommand{\ctLowHalf}{202}
\newcommand{\cvK}{10}
\newcommand{\cvLassoErrCv}{0.724}
\newcommand{\cvLassoErrOracle}{0.624}
\newcommand{\cvLassoExtra}{16}
\newcommand{\cvLassoHi}{0.316}
\newcommand{\cvLassoLo}{0.123}
\newcommand{\cvLassoMed}{0.200}
\newcommand{\cvLassoOracle}{0.200}
\newcommand{\cvReps}{100}
\newcommand{\cvRidgeErrCv}{0.997}
\newcommand{\cvRidgeErrOracle}{0.915}
\newcommand{\cvRidgeExtra}{9}
\newcommand{\cvRidgeHi}{177.8}
\newcommand{\cvRidgeLo}{31.6}
\newcommand{\cvRidgeMed}{56.2}
\newcommand{\cvRidgeOracle}{75.0}
\newcommand{\dpFrom}{1926}
\newcommand{\dpHalf}{9}
\newcommand{\dpMoveMed}{6}
\newcommand{\dpRho}{0.9938}
\newcommand{\dpSimHalf}{3}
\newcommand{\dpSimRho}{0.969}
\newcommand{\dpTo}{2025}
\newcommand{\dsAtModels}{0.718}
\newcommand{\dsCross}{5}
\newcommand{\dsHundred}{0.408}
\newcommand{\dsTwo}{0.993}
\newcommand{\eqBestUni}{Treasury bill rate}
\newcommand{\eqBestUniRtwo}{0.08}
\newcommand{\eqBoostM}{150}
\newcommand{\eqBoostNu}{0.02}
\newcommand{\eqCeBm}{\ensuremath{-}1.98}
\newcommand{\eqCeBoost}{\ensuremath{-}1.76}
\newcommand{\eqCeComb}{0.48}
\newcommand{\eqCeDe}{\ensuremath{-}0.51}
\newcommand{\eqCeDfr}{0.80}
\newcommand{\eqCeDfy}{\ensuremath{-}0.09}
\newcommand{\eqCeDp}{\ensuremath{-}1.03}
\newcommand{\eqCeDy}{\ensuremath{-}0.70}
\newcommand{\eqCeEp}{\ensuremath{-}0.30}
\newcommand{\eqCeForest}{\ensuremath{-}1.25}
\newcommand{\eqCeHm}{6.49}
\newcommand{\eqCeInfl}{0.15}
\newcommand{\eqCeLasso}{0.22}
\newcommand{\eqCeLtr}{0.29}
\newcommand{\eqCeLty}{1.08}
\newcommand{\eqCeNtis}{0.91}
\newcommand{\eqCeOls}{0.80}
\newcommand{\eqCeOlsct}{0.80}
\newcommand{\eqCeRidge}{1.01}
\newcommand{\eqCeSvar}{0.44}
\newcommand{\eqCeTbl}{1.04}
\newcommand{\eqCeTms}{1.06}
\newcommand{\eqCtComb}{28}
\newcommand{\eqCtHalf}{51}
\newcommand{\eqCwBm}{0.09}
\newcommand{\eqCwBoost}{\ensuremath{-}1.44}
\newcommand{\eqCwComb}{1.51}
\newcommand{\eqCwDe}{\ensuremath{-}1.71}
\newcommand{\eqCwDfr}{0.08}
\newcommand{\eqCwDfy}{\ensuremath{-}0.27}
\newcommand{\eqCwDp}{0.81}
\newcommand{\eqCwDy}{0.97}
\newcommand{\eqCwEp}{0.37}
\newcommand{\eqCwForest}{\ensuremath{-}0.31}
\newcommand{\eqCwInfl}{0.23}
\newcommand{\eqCwLasso}{\ensuremath{-}0.88}
\newcommand{\eqCwLtr}{0.56}
\newcommand{\eqCwLty}{1.28}
\newcommand{\eqCwNtis}{0.48}
\newcommand{\eqCwOls}{0.08}
\newcommand{\eqCwOlsct}{0.27}
\newcommand{\eqCwRidge}{\ensuremath{-}0.22}
\newcommand{\eqCwSvar}{\ensuremath{-}0.47}
\newcommand{\eqCwTbl}{1.42}
\newcommand{\eqCwTms}{0.84}
\newcommand{\eqFirst}{1926}
\newcommand{\eqForestB}{100}
\newcommand{\eqForestMtry}{5}
\newcommand{\eqGamma}{3}
\newcommand{\eqLassoNz}{1.3}
\newcommand{\eqLeaf}{40}
\newcommand{\eqOlsGainDecade}{2010s}
\newcommand{\eqOlsSecondYear}{2020}
\newcommand{\eqOlsWorstYear}{2009}
\newcommand{\eqOosFrom}{1965}
\newcommand{\eqOosN}{732}
\newcommand{\eqOosTo}{2025}
\newcommand{\eqPBm}{0.465}
\newcommand{\eqPBoost}{0.925}
\newcommand{\eqPComb}{0.066}
\newcommand{\eqPDe}{0.957}
\newcommand{\eqPDfr}{0.466}
\newcommand{\eqPDfy}{0.605}
\newcommand{\eqPDp}{0.209}
\newcommand{\eqPDy}{0.166}
\newcommand{\eqPEp}{0.355}
\newcommand{\eqPForest}{0.620}
\newcommand{\eqPInfl}{0.410}
\newcommand{\eqPLasso}{0.811}
\newcommand{\eqPLtr}{0.289}
\newcommand{\eqPLty}{0.100}
\newcommand{\eqPNtis}{0.315}
\newcommand{\eqPOls}{0.470}
\newcommand{\eqPOlsct}{0.392}
\newcommand{\eqPRidge}{0.587}
\newcommand{\eqPSvar}{0.679}
\newcommand{\eqPTbl}{0.078}
\newcommand{\eqPTms}{0.200}
\newcommand{\eqPred}{14}
\newcommand{\eqRtwoABm}{\ensuremath{-}3.29}
\newcommand{\eqRtwoABoost}{\ensuremath{-}3.97}
\newcommand{\eqRtwoAComb}{0.44}
\newcommand{\eqRtwoADe}{\ensuremath{-}1.70}
\newcommand{\eqRtwoADfr}{0.07}
\newcommand{\eqRtwoADfy}{0.03}
\newcommand{\eqRtwoADp}{\ensuremath{-}0.53}
\newcommand{\eqRtwoADy}{\ensuremath{-}0.93}
\newcommand{\eqRtwoAEp}{\ensuremath{-}2.18}
\newcommand{\eqRtwoAForest}{\ensuremath{-}1.38}
\newcommand{\eqRtwoAInfl}{0.28}
\newcommand{\eqRtwoALasso}{\ensuremath{-}0.02}
\newcommand{\eqRtwoALtr}{\ensuremath{-}0.17}
\newcommand{\eqRtwoALty}{\ensuremath{-}1.20}
\newcommand{\eqRtwoANtis}{0.49}
\newcommand{\eqRtwoAOls}{\ensuremath{-}7.51}
\newcommand{\eqRtwoAOlsct}{\ensuremath{-}1.47}
\newcommand{\eqRtwoARidge}{0.57}
\newcommand{\eqRtwoASvar}{\ensuremath{-}0.14}
\newcommand{\eqRtwoATbl}{\ensuremath{-}0.04}
\newcommand{\eqRtwoATms}{0.60}
\newcommand{\eqRtwoBBm}{\ensuremath{-}0.17}
\newcommand{\eqRtwoBBoost}{\ensuremath{-}7.34}
\newcommand{\eqRtwoBComb}{0.07}
\newcommand{\eqRtwoBDe}{\ensuremath{-}1.33}
\newcommand{\eqRtwoBDfr}{\ensuremath{-}0.63}
\newcommand{\eqRtwoBDfy}{\ensuremath{-}0.12}
\newcommand{\eqRtwoBDp}{\ensuremath{-}0.08}
\newcommand{\eqRtwoBDy}{\ensuremath{-}0.16}
\newcommand{\eqRtwoBEp}{\ensuremath{-}0.82}
\newcommand{\eqRtwoBForest}{\ensuremath{-}4.27}
\newcommand{\eqRtwoBInfl}{\ensuremath{-}0.42}
\newcommand{\eqRtwoBLasso}{\ensuremath{-}2.94}
\newcommand{\eqRtwoBLtr}{\ensuremath{-}0.35}
\newcommand{\eqRtwoBLty}{0.48}
\newcommand{\eqRtwoBNtis}{\ensuremath{-}2.18}
\newcommand{\eqRtwoBOls}{\ensuremath{-}8.43}
\newcommand{\eqRtwoBOlsct}{\ensuremath{-}2.76}
\newcommand{\eqRtwoBRidge}{\ensuremath{-}2.67}
\newcommand{\eqRtwoBSvar}{\ensuremath{-}0.61}
\newcommand{\eqRtwoBTbl}{0.24}
\newcommand{\eqRtwoBTms}{\ensuremath{-}0.94}
\newcommand{\eqRtwoBm}{\ensuremath{-}1.93}
\newcommand{\eqRtwoBoost}{\ensuremath{-}5.43}
\newcommand{\eqRtwoComb}{0.28}
\newcommand{\eqRtwoDe}{\ensuremath{-}1.54}
\newcommand{\eqRtwoDfr}{\ensuremath{-}0.24}
\newcommand{\eqRtwoDfy}{\ensuremath{-}0.04}
\newcommand{\eqRtwoDp}{\ensuremath{-}0.33}
\newcommand{\eqRtwoDy}{\ensuremath{-}0.60}
\newcommand{\eqRtwoEp}{\ensuremath{-}1.59}
\newcommand{\eqRtwoForest}{\ensuremath{-}2.63}
\newcommand{\eqRtwoInfl}{\ensuremath{-}0.02}
\newcommand{\eqRtwoLasso}{\ensuremath{-}1.29}
\newcommand{\eqRtwoLtr}{\ensuremath{-}0.25}
\newcommand{\eqRtwoLty}{\ensuremath{-}0.47}
\newcommand{\eqRtwoNtis}{\ensuremath{-}0.67}
\newcommand{\eqRtwoOls}{\ensuremath{-}7.91}
\newcommand{\eqRtwoOlsct}{\ensuremath{-}2.03}
\newcommand{\eqRtwoRidge}{\ensuremath{-}0.83}
\newcommand{\eqRtwoSvar}{\ensuremath{-}0.35}
\newcommand{\eqRtwoTbl}{0.08}
\newcommand{\eqRtwoTms}{\ensuremath{-}0.07}
\newcommand{\eqS}{0.100}
\newcommand{\eqSAnn}{0.35}
\newcommand{\eqSrBm}{0.19}
\newcommand{\eqSrBoost}{0.25}
\newcommand{\eqSrComb}{0.39}
\newcommand{\eqSrDe}{0.36}
\newcommand{\eqSrDfr}{0.42}
\newcommand{\eqSrDfy}{0.36}
\newcommand{\eqSrDp}{0.26}
\newcommand{\eqSrDy}{0.29}
\newcommand{\eqSrEp}{0.33}
\newcommand{\eqSrForest}{0.28}
\newcommand{\eqSrHm}{0.37}
\newcommand{\eqSrInfl}{0.37}
\newcommand{\eqSrLasso}{0.38}
\newcommand{\eqSrLtr}{0.39}
\newcommand{\eqSrLty}{0.44}
\newcommand{\eqSrNtis}{0.44}
\newcommand{\eqSrOls}{0.42}
\newcommand{\eqSrOlsct}{0.42}
\newcommand{\eqSrRidge}{0.43}
\newcommand{\eqSrSvar}{0.39}
\newcommand{\eqSrTbl}{0.43}
\newcommand{\eqSrTms}{0.44}
\newcommand{\eqUniPos}{1}
\newcommand{\eqUniSig}{0}
\newcommand{\eqValid}{120}
\newcommand{\eqVarWin}{60}
\newcommand{\eqWhi}{1.5}
\newcommand{\eqWinFirst}{456}
\newcommand{\eqWinLast}{1,187}
\newcommand{\eqWorstUni}{book to market}
\newcommand{\eqWorstUniRtwo}{\ensuremath{-}1.9}
\newcommand{\gmEdgeHigh}{0.60}
\newcommand{\gmEdgeLow}{0.24}
\newcommand{\gmErrHigh}{0.21}
\newcommand{\gmErrLow}{0.32}
\newcommand{\gmNegHigh}{4}
\newcommand{\gmNegLow}{7}
\newcommand{\gmPeakHigh}{1.17}
\newcommand{\gmPeakLow}{4.03}
\newcommand{\hdBandFlat}{0.044}
\newcommand{\hdBandWw}{0.57}
\newcommand{\hdBs}{3.988}
\newcommand{\hdBumps}{11}
\newcommand{\hdCStdBs}{0.503}
\newcommand{\hdCStdCost}{0.497}
\newcommand{\hdCStdFlat}{0.604}
\newcommand{\hdCStdWw}{0.619}
\newcommand{\hdCStdWwOne}{0.784}
\newcommand{\hdCStdZero}{0.502}
\newcommand{\hdCentre}{0.5}
\newcommand{\hdCostBs}{0.709}
\newcommand{\hdCostCost}{0.642}
\newcommand{\hdCostFlat}{0.448}
\newcommand{\hdCostWw}{0.434}
\newcommand{\hdCostWwOne}{0.368}
\newcommand{\hdCostZero}{0.710}
\newcommand{\hdDeltaGap}{0.002}
\newcommand{\hdGamma}{1}
\newcommand{\hdIters}{500}
\newcommand{\hdK}{100}
\newcommand{\hdKappa}{0.2}
\newcommand{\hdMarkupBs}{0.832}
\newcommand{\hdMarkupCost}{0.763}
\newcommand{\hdMarkupFlat}{0.628}
\newcommand{\hdMarkupWw}{0.623}
\newcommand{\hdMarkupWwOne}{0.673}
\newcommand{\hdMarkupZero}{0.833}
\newcommand{\hdMeanBs}{0.003}
\newcommand{\hdMeanLearn}{0.006}
\newcommand{\hdNte}{20,000}
\newcommand{\hdNtr}{20,000}
\newcommand{\hdParams}{23}
\newcommand{\hdPriceBs}{4.820}
\newcommand{\hdPriceCost}{4.751}
\newcommand{\hdPriceCutCost}{8}
\newcommand{\hdPriceCutWw}{25}
\newcommand{\hdPriceFlat}{4.616}
\newcommand{\hdPriceFree}{4.079}
\newcommand{\hdPriceLearn}{3.990}
\newcommand{\hdPriceWw}{4.610}
\newcommand{\hdPriceWwOne}{4.661}
\newcommand{\hdPriceZero}{4.821}
\newcommand{\hdSZero}{100}
\newcommand{\hdSig}{20}
\newcommand{\hdStdBs}{0.434}
\newcommand{\hdStdLearn}{0.434}
\newcommand{\hdStdLearnTrain}{0.438}
\newcommand{\hdSteps}{63}
\newcommand{\hdTurnBs}{3.51}
\newcommand{\hdTurnCost}{3.18}
\newcommand{\hdTurnFlat}{2.20}
\newcommand{\hdTurnWw}{2.13}
\newcommand{\hdTurnWwOne}{1.80}
\newcommand{\hdTurnZero}{3.51}
\newcommand{\hdUnhedged}{6.14}
\newcommand{\hdWidth}{0.100}
\newcommand{\hdWwAtm}{0.078}
\newcommand{\hpDays}{63}
\newcommand{\hpTrades}{32}
\newcommand{\hpTurnBand}{1.09}
\newcommand{\hpTurnDelta}{1.95}
\newcommand{\kfBlockMean}{\ensuremath{-}2.6}
\newcommand{\kfBlockMed}{\ensuremath{-}2.8}
\newcommand{\kfBlockPos}{21}
\newcommand{\kfCorrEleven}{0.088}
\newcommand{\kfCorrOne}{0.917}
\newcommand{\kfCorrSix}{0.501}
\newcommand{\kfCorrThEleven}{0.083}
\newcommand{\kfCorrThOne}{0.917}
\newcommand{\kfCorrThSix}{0.500}
\newcommand{\kfEmb}{12}
\newcommand{\kfFreshMean}{\ensuremath{-}3.5}
\newcommand{\kfFreshMed}{\ensuremath{-}1.0}
\newcommand{\kfFreshPos}{35}
\newcommand{\kfH}{12}
\newcommand{\kfK}{10}
\newcommand{\kfMu}{0.5}
\newcommand{\kfN}{468}
\newcommand{\kfPhi}{0.98}
\newcommand{\kfPurgeMean}{\ensuremath{-}3.6}
\newcommand{\kfPurgeMed}{\ensuremath{-}3.6}
\newcommand{\kfPurgePos}{21}
\newcommand{\kfReps}{1,000}
\newcommand{\kfSelIn}{\ensuremath{-}1.97}
\newcommand{\kfSelInPos}{16}
\newcommand{\kfSelN}{100}
\newcommand{\kfSelOut}{1.09}
\newcommand{\kfSelOutPos}{98}
\newcommand{\kfSelReps}{400}
\newcommand{\kfSelT}{480}
\newcommand{\kfShufHigh}{12}
\newcommand{\kfShufMean}{1.8}
\newcommand{\kfShufMed}{0.7}
\newcommand{\kfShufPos}{67}
\newcommand{\kfT}{480}
\newcommand{\kfVol}{4.5}
\newcommand{\mtAssets}{6}
\newcommand{\mtB}{5,000}
\newcommand{\mtBestAsset}{WTI crude}
\newcommand{\mtBestLook}{3}
\newcommand{\mtBestMean}{10.99}
\newcommand{\mtBestP}{0.058}
\newcommand{\mtBestT}{1.61}
\newcommand{\mtBhFive}{0}
\newcommand{\mtBhTen}{0}
\newcommand{\mtBhyFive}{0}
\newcommand{\mtBhyTen}{0}
\newcommand{\mtBlock}{6}
\newcommand{\mtCm}{4.77}
\newcommand{\mtCost}{0.1}
\newcommand{\mtCrit}{2.95}
\newcommand{\mtFirst}{March 2002}
\newcommand{\mtHolmFive}{0}
\newcommand{\mtHolmTen}{0}
\newcommand{\mtLast}{September 2026}
\newcommand{\mtLookMax}{12}
\newcommand{\mtLookMin}{2}
\newcommand{\mtN}{295}
\newcommand{\mtNaiveFive}{4}
\newcommand{\mtNaiveTen}{11}
\newcommand{\mtPos}{30}
\newcommand{\mtRcP}{0.178}
\newcommand{\mtRules}{66}
\newcommand{\mtSimB}{499}
\newcommand{\mtSimBhy}{7.4}
\newcommand{\mtSimNaive}{62.2}
\newcommand{\mtSimRc}{3.3}
\newcommand{\mtSimReps}{1000}
\newcommand{\mtSimSpa}{6.9}
\newcommand{\mtSimStepM}{7.0}
\newcommand{\mtSpaP}{0.181}
\newcommand{\mtSpaPl}{0.120}
\newcommand{\mtSpaPu}{0.185}
\newcommand{\mtStepM}{0}
\newcommand{\mtThr}{\ensuremath{-}1.86}
\newcommand{\mtTopAsset}{euro}
\newcommand{\mtTopLook}{4}
\newcommand{\mtTopP}{0.011}
\newcommand{\mtTopT}{2.39}
\newcommand{\mtTurn}{3.1}
\newcommand{\nnBatch}{128}
\newcommand{\nnConvex}{4.8}
\newcommand{\nnDeltaMax}{0.061}
\newcommand{\nnDeltaMaxM}{1.25}
\newcommand{\nnDeltaMaxV}{0.50}
\newcommand{\nnDeltaRmse}{0.0070}
\newcommand{\nnEpochs}{4,000}
\newcommand{\nnHid}{32}
\newcommand{\nnLr}{0.005}
\newcommand{\nnLrEnd}{0.00005}
\newcommand{\nnMax}{21.1}
\newcommand{\nnMhi}{1.25}
\newcommand{\nnMlo}{0.75}
\newcommand{\nnNtr}{4,000}
\newcommand{\nnParams}{1,185}
\newcommand{\nnPolyDeg}{6}
\newcommand{\nnPolyMax}{38.3}
\newcommand{\nnPolyRmse}{6.0}
\newcommand{\nnPolyTerms}{28}
\newcommand{\nnRmse}{2.2}
\newcommand{\nnSteps}{128,000}
\newcommand{\nnTrainLoss}{3.73}
\newcommand{\nnValLoss}{4.58}
\newcommand{\nnVhi}{0.5}
\newcommand{\nnVlo}{0.05}
\newcommand{\ovApproxHundred}{2.531}
\newcommand{\ovApproxTen}{1.575}
\newcommand{\ovApproxThousand}{3.255}
\newcommand{\ovBestIs}{0.67}
\newcommand{\ovBestOos}{0.00}
\newcommand{\ovBestOosSe}{0.02}
\newcommand{\ovBhFive}{0}
\newcommand{\ovBhTen}{0}
\newcommand{\ovBonfFive}{0}
\newcommand{\ovBonfTen}{0}
\newcommand{\ovBoundHundred}{3.035}
\newcommand{\ovBoundTen}{2.146}
\newcommand{\ovBoundThousand}{3.717}
\newcommand{\ovCorr}{0.52}
\newcommand{\ovEqBench}{0.0932}
\newcommand{\ovEqBest}{net equity expansion}
\newcommand{\ovEqDistinct}{20}
\newcommand{\ovEqDsr}{0.718}
\newcommand{\ovEqDsrDistinct}{0.706}
\newcommand{\ovEqKurt}{6.7}
\newcommand{\ovEqModels}{21}
\newcommand{\ovEqPsr}{0.999}
\newcommand{\ovEqSe}{0.0364}
\newcommand{\ovEqSkew}{0.53}
\newcommand{\ovEqSr}{0.1142}
\newcommand{\ovEqSrAnn}{0.40}
\newcommand{\ovEqSrSd}{0.0485}
\newcommand{\ovEqT}{732}
\newcommand{\ovHalf}{94}
\newcommand{\ovHolmFirst}{0.0071}
\newcommand{\ovHolmFive}{0}
\newcommand{\ovHolmTen}{0}
\newcommand{\ovLook}{50}
\newcommand{\ovMaxHundred}{2.508}
\newcommand{\ovMaxSrHundred}{0.79}
\newcommand{\ovMaxTen}{1.539}
\newcommand{\ovMaxThousand}{3.241}
\newcommand{\ovNeff}{20}
\newcommand{\ovPbo}{0.55}
\newcommand{\ovPboFig}{0.39}
\newcommand{\ovPboRuns}{50}
\newcommand{\ovPmin}{0.078}
\newcommand{\ovRawFive}{0}
\newcommand{\ovRawTen}{1}
\newcommand{\ovReps}{200}
\newcommand{\ovRules}{200}
\newcommand{\ovSFig}{16}
\newcommand{\ovSRep}{10}
\newcommand{\ovSimHundred}{2.508}
\newcommand{\ovSimReps}{20,000}
\newcommand{\ovSimTen}{1.539}
\newcommand{\ovSimThousand}{3.245}
\newcommand{\ovSlope}{\ensuremath{-}0.49}
\newcommand{\ovSplits}{12,870}
\newcommand{\ovT}{240}
\newcommand{\ovVol}{4.5}
\newcommand{\ovYears}{10}
\newcommand{\skHalfBig}{8}
\newcommand{\skHalfSmall}{139}
\newcommand{\tbB}{300}
\newcommand{\tbBag}{3.80}
\newcommand{\tbBoost}{1.76}
\newcommand{\tbBoostAt}{278}
\newcommand{\tbBoostBest}{1.73}
\newcommand{\tbBoostHundred}{1.90}
\newcommand{\tbBoostOne}{5.26}
\newcommand{\tbBoostOneAt}{9}
\newcommand{\tbBoostOneEnd}{6.13}
\newcommand{\tbDeep}{8.81}
\newcommand{\tbDeepLeaves}{81}
\newcommand{\tbDepth}{9}
\newcommand{\tbM}{1,000}
\newcommand{\tbMtry}{3}
\newcommand{\tbN}{500}
\newcommand{\tbNte}{5,000}
\newcommand{\tbNu}{0.1}
\newcommand{\tbOls}{6.14}
\newcommand{\tbP}{10}
\newcommand{\tbRf}{4.40}
\newcommand{\tbRho}{0.038}
\newcommand{\tbSigTwo}{9.54}
\newcommand{\tbTree}{8.66}
\newcommand{\tbTreeLeaves}{76}
\newcommand{\tbVarAvg}{0.56}
\newcommand{\tbVarB}{50}
\newcommand{\tbVarF}{24.4}
\newcommand{\tbVarFPop}{23.8}
\newcommand{\tbVarFormula}{0.55}
\newcommand{\tbVarR}{40}
\newcommand{\zaLasso}{0.55}
\newcommand{\zaLassoSmall}{0}
\newcommand{\zaRidge}{0.64}
\newcommand{\zbC}{0.667}
\newcommand{\zbMse}{0.333}
\newcommand{\zcVar}{0.307}
\newcommand{\zcVarTen}{0.37}
\newcommand{\zdCorr}{0.667}
\newcommand{\zdPurged}{28}
\newcommand{\zeWone}{1.4537}
\newcommand{\zeWtwo}{0.4271}
\newcommand{\zeYhat}{0.9242}
\newcommand{\zfRtwo}{1.04}
\newcommand{\zgGain}{35.4}
\newcommand{\zhAdj}{0.00006}
\newcommand{\zhGap}{0.00002}
\newcommand{\ziExact}{2.508}
\newcommand{\ziMax}{2.531}
\newcommand{\ziSr}{1.13}
\newcommand{\zjBh}{4}
\newcommand{\zjBonf}{2}
\newcommand{\zjHolm}{3}
\newcommand{\zkPsr}{0.9941}
\newcommand{\zkSe}{0.0992}
\newcommand{\zkZ}{2.519}
\newcommand{\zlCut}{84.4}
\newcommand{\zlH}{0.5625}

\begin{chapterrefs}
\setcounter{chapter}{11}
\chapter{Machine Learning in Finance}
\label{ch:ml}

Machine learning is a set of methods for fitting flexible functions to data and for judging those
functions by how well they predict data they have not seen. Finance offers it a hard test. Returns are
noisy, the signal in them is small, observations arrive in time order and overlap, and the same history
has been searched by many people before. A method that finds structure in a sample of images or
sentences will usually find structure in a sample of returns too, whether or not the structure is
there. This chapter is about telling the two apart.

It takes six questions in turn. Why does a biased estimator often predict better than an unbiased one,
and how do ridge regression and the lasso decide how much bias to take on? How should a model be validated when the
observations are ordered in time and their labels overlap? What do trees, forests and boosting add to
a linear model, and what limits each of them? How is a neural network trained, and can such models learn
a price, its delta and a hedge when trading is costly? Can any of these methods forecast the equity premium out of
sample, and what would a small forecasting gain be worth to an investor? And how much of a good
backtest is the search that produced it?

Several earlier chapters are used without being repeated. The standard error of a Sharpe ratio, with
its corrections for skewness and kurtosis, is derived in Chapter~8, and screening many managers with
false discovery control is in Chapter~7; \cref{sec12:overfit} builds the deflated Sharpe ratio,
Holm's procedure and the bootstrap tests of data snooping on top of them. The discrete hedging error of a Black-Scholes delta is in Chapter~5,
and \cref{sec12:hedge} uses it as the benchmark a learned hedge has to match.

\begin{keybox}[How to read this chapter]
The numbers in the text are computed by nine scripts in \texttt{code/ch12/} (\texttt{regularise},
\texttt{cv\_time}, \texttt{trees}, \texttt{neural}, \texttt{equity\_premium}, \texttt{hedging},
\texttt{overfit}, \texttt{snooping} and \texttt{exercises}). Every method is written out in the chapter's library
\texttt{ch12lib.py} with NumPy and SciPy: ridge through the singular value decomposition, the lasso by
coordinate descent, regression trees, bagging, forests and boosting, a feed-forward network with its
backward pass and the Adam optimiser, the out-of-sample statistics, the learned hedges and the
overfitting statistics. The tests in \texttt{code/ch12/test\_*.py} check the library against closed
forms, finite differences and identities, and compare the lasso and the regression tree with
scikit-learn \parencite{sklearn2011}, which is used for that comparison only; ridge is checked against
the normal equations. An independent R script
rebuilds the fourteen predictors from the data file and recomputes every least squares forecast, the
out-of-sample $R^2$, the Clark and West statistics, the certainty equivalents, the deflated Sharpe ratio,
the expected maximum of normal variables and the multiple testing counts. GNU Octave recomputes the first
ridge forecast by the normal equations, the first lasso forecast by proximal gradient steps, the
Black-Scholes price by quadrature and the exercise answers.

The equity premium data are the monthly predictors of Welch and Goyal, updated to December 2025
\parencite{welch2008,goyalch12}. The file is free to download but carries no licence for
redistribution, so it is not shipped; \texttt{scripts/pull\_ch12.py} downloads it and builds the chapter's data file. Every
other data set in the chapter is simulated, and the simulations are labelled as such.

Random draws come from NumPy's default generator with base seed \texttt{8317} and a stream number fixed
in each script. The seed is arbitrary; fixing it makes every run repeatable. Each short listing that
draws random numbers fixes its own seed, distinct from \texttt{8317}, so that it runs on its own.

Notation. $X$ is an $n\times p$ matrix of predictors with rows $x_i^\top$, $y$ a vector of responses,
$\beta$ a coefficient vector and $\varepsilon$ a noise term with variance $\sigma^2$. In
\cref{sec12:reg} the predictors and the response are centred, so no intercept is needed. $\lambda$ is a
penalty, $\|\cdot\|_2$ and $\|\cdot\|_1$ are the Euclidean and absolute value norms, and, when $p\le n$ and $X$ has
full column rank, $X=UDV^\top$ is a thin singular value decomposition with singular values
$d_1\ge\dots\ge d_p>0$; otherwise the decomposition is taken over the $\operatorname{rank}(X)$ nonzero
singular values. In \cref{sec12:cv}, $h$ is the
number of periods a label spans. In \cref{sec12:nn}, $W_\ell$ and $b_\ell$ are the weights and biases of
layer $\ell$, $a_\ell$ its output, $J$ the loss, and $\delta_\ell$ the gradient of $J$ with respect to the layer's
input $z_\ell$. In \cref{sec12:premium}, $r_{t+1}$ is the log equity premium of month $t+1$, $\bar r_t$ the
historical mean through month $t$ and $\hat r_{t+1}$ a model's forecast. In \cref{sec12:hedge}, $S_k$ is
the stock price at date $k$, $\Delta_k$ the hedge ratio held from $k$ to $k+1$, $\kappa$ the proportional
cost, $\gamma$ a risk aversion and $\mathcal L$ the hedger's loss. In \cref{sec12:overfit}, $\hat\theta$ is a per-period Sharpe
ratio estimated from $T$ observations, $\theta^*$ a benchmark for it and $N$ a number of trials.
\end{keybox}

\FloatBarrier
% =====================================================================================
\section{Why can a biased estimator predict better?}
\label{sec12:reg}

A forecast is judged on new data. Suppose $y=f(x)+\varepsilon$ with $\E\varepsilon=0$ and
$\operatorname{Var}\varepsilon=\sigma^2$, and a method fitted to a training sample\ix{Training set} $\mathcal D$ returns a
function $\hat f$. The training sample is random, so $\hat f(x_0)$ at a fixed point $x_0$ is a random
variable, and its error splits into three parts.

\begin{proposition}[Bias and variance]\label{prop12:bv}
Let $y_0=f(x_0)+\varepsilon_0$ be a new response, with $\varepsilon_0$ independent of $\mathcal D$. Then
\begin{equation}\label{eq12:bv}
  \E\bigl[(y_0-\hat f(x_0))^2\bigr]=\sigma^2+\bigl(\E\hat f(x_0)-f(x_0)\bigr)^2+\operatorname{Var}\hat f(x_0),
\end{equation}
where the expectations are over $\mathcal D$ and $\varepsilon_0$.
\end{proposition}

\begin{proof}
Write $y_0-\hat f(x_0)=\varepsilon_0+\bigl(f(x_0)-\E\hat f(x_0)\bigr)+\bigl(\E\hat f(x_0)-\hat f(x_0)\bigr)$. The
middle term is a constant, the last has mean zero, and $\varepsilon_0$ has mean zero and is independent
of the other two. Squaring and taking expectations, every cross product vanishes.
\end{proof}

The noise $\sigma^2$ is beyond any method. The other two terms trade against each other. Least
squares on the correct linear model is unbiased, so all of its excess error is variance, and with many
predictors and few observations that variance is large. A method that shrinks the coefficients towards
zero adds bias\ix{Bias} and removes variance, and the sum can fall.

\subsection{Ridge regression}

Ridge\ix{Ridge regression} regression\ix{Linear regression} \parencite{hoerl1970} adds a penalty on the squared size of the coefficients:
\begin{equation}\label{eq12:ridge}
  \hat\beta_\lambda=\arg\min_\beta\ \|y-X\beta\|_2^2+\lambda\|\beta\|_2^2=(X^\top X+\lambda I)^{-1}X^\top y .
\end{equation}
The inverse exists for every $\lambda>0$, even when $p>n$ and least squares\ix{Least squares} has no unique solution.

\begin{proposition}[What ridge does]\label{prop12:ridge}
Let $X=UDV^\top$ and $y=X\beta+\varepsilon$ with $\E\varepsilon=0$ and $\operatorname{Var}\varepsilon=\sigma^2I$,
for a fixed $X$ of rank $p$, and write $\alpha=V^\top\beta$.
\begin{enumerate}[label=(\roman*)]
\item The fitted values are $X\hat\beta_\lambda=\sum_j u_j\,\dfrac{d_j^2}{d_j^2+\lambda}\,u_j^\top y$: each principal
direction of the predictors is shrunk by the factor $d_j^2/(d_j^2+\lambda)$, and most where $d_j$ is small.
\item The mean squared error of the coefficients is
\begin{equation}\label{eq12:ridgemse}
  M(\lambda)=\E\|\hat\beta_\lambda-\beta\|_2^2=\sum_{j=1}^p\frac{\sigma^2d_j^2+\lambda^2\alpha_j^2}{(d_j^2+\lambda)^2}.
\end{equation}
\item If $\sigma^2>0$, then $M'(0)=-2\sigma^2\sum_jd_j^{-4}<0$, so a small enough $\lambda>0$ gives a smaller
error than least squares, whatever $\beta$ is.
\end{enumerate}
\end{proposition}

\begin{proof}
(i) With $X^\top X=VD^2V^\top$, $\hat\beta_\lambda=V(D^2+\lambda I)^{-1}DU^\top y$, and multiplying by
$X=UDV^\top$ gives the sum. (ii) In the rotated coordinates $\hat\alpha=V^\top\hat\beta_\lambda$ has
components $\hat\alpha_j=d_j(u_j^\top y)/(d_j^2+\lambda)$, and $u_j^\top y=d_j\alpha_j+u_j^\top\varepsilon$
with $\operatorname{Var}(u_j^\top\varepsilon)=\sigma^2$. So $\E\hat\alpha_j-\alpha_j=-\lambda\alpha_j/(d_j^2+\lambda)$
and $\operatorname{Var}\hat\alpha_j=\sigma^2d_j^2/(d_j^2+\lambda)^2$. The rotation preserves length, and the sum
of squared bias and variance\ix{Bias-variance trade-off}\index{Bias-variance trade-off} over $j$ is \cref{eq12:ridgemse}. (iii) The derivative of the $j$th term is
\[
  \frac{2\lambda\alpha_j^2(d_j^2+\lambda)-2(\sigma^2d_j^2+\lambda^2\alpha_j^2)}{(d_j^2+\lambda)^3},
\]
which at $\lambda=0$ is $-2\sigma^2/d_j^4$. The derivative is continuous, so $M(\lambda)<M(0)$ for all small
enough $\lambda>0$.
\end{proof}

The listing below computes ridge from the thin singular value decomposition\ix{Singular value decomposition} of part (i), without
forming an inverse. It is written for a predictor matrix of full column rank. Singular values below a
relative tolerance are treated as zero and given a factor of zero, so a rank-deficient matrix at
$\lambda=0$ gives the minimum-norm solution instead of dividing by a rounding error. It returns the factor $d_j^2/(d_j^2+\lambda)$ by which each singular direction is
shrunk and their sum, the effective degrees of freedom\ix{Degrees of freedom}\ix{Effective degrees of freedom}, and checks the coefficients against the normal
equations\ix{Normal equations} on a small simulated sample: with six predictors the degrees of freedom fall from 6 at
$\lambda=0$ to 4.76 at $\lambda=25$ and 1.72 at $\lambda=250$.
\Needspace*{22\baselineskip}
\lstinputlisting[language=Python]{code/ch12/snippets/ridge_svd.py}

Part (iii) is the existence result of \textcite{hoerl1970}. It does not say which $\lambda$ to use, and
the best one depends on the unknown $\alpha$; in practice $\lambda$ is chosen by cross-validation\index{Cross-validation}. The
trace of the hat matrix, $\sum_jd_j^2/(d_j^2+\lambda)$, falls from $p$ to zero as $\lambda$ grows and is
called the effective degrees of freedom of the fit.

\subsection{The lasso}

The lasso\ix{Lasso} \parencite{tibshirani1996} replaces the squared penalty by the absolute value:
\begin{equation}\label{eq12:lasso}
  \hat\beta=\arg\min_\beta\ \frac1{2n}\|y-X\beta\|_2^2+\lambda\|\beta\|_1 .
\end{equation}
The objective is convex but not differentiable where a coefficient is zero, and that kink is what sets
coefficients exactly to zero. The factor $1/(2n)$ follows \textcite{friedman2010} and puts $\lambda$ on the scale of a coefficient; the
$\lambda$ of ridge in \cref{eq12:ridge} is on the scale of $n$, so the two penalties quoted below are not
comparable numbers. Let $S(z,t)=\operatorname{sign}(z)\max(|z|-t,0)$ be the soft-threshold function.

\begin{proposition}[Optimality conditions of the lasso]\label{prop12:lasso}
\begin{enumerate}[label=(\roman*)]
\item $\hat\beta$ solves \cref{eq12:lasso} if and only if, with $g=X^\top(y-X\hat\beta)/n$,
$g_j=\lambda\operatorname{sign}\hat\beta_j$ for every $j$ with $\hat\beta_j\ne0$ and $|g_j|\le\lambda$ for every $j$
with $\hat\beta_j=0$.
\item $\hat\beta=0$ is a solution if and only if $\lambda\ge\lambda_{\max}=\max_j|x_{(j)}^\top y|/n$, where $x_{(j)}$ is the
$j$th column of $X$.
\item If $X^\top X=nI$, then $\hat\beta_j=S(x_{(j)}^\top y/n,\lambda)$: each least squares coefficient is
moved towards zero by $\lambda$ and set to zero if its absolute value is at most $\lambda$.
\item With all coordinates but the $j$th held fixed, the objective is minimised over $\beta_j$ by
$S(z_j,\lambda)/c_j$, where $c_j=\|x_{(j)}\|_2^2/n$ and $z_j=x_{(j)}^\top r/n+c_j\beta_j$ for the current
residual $r=y-X\beta$.
\end{enumerate}
\end{proposition}

\begin{proof}
(i) A convex function is minimised exactly where zero belongs to its subdifferential. The
subdifferential of $|\beta_j|$ is $\{\operatorname{sign}\beta_j\}$ when $\beta_j\ne0$ and the interval $[-1,1]$
when $\beta_j=0$, and the gradient of the quadratic part is $-g$; so the condition is $g_j\in\lambda\,\partial|\hat\beta_j|$
for every $j$, which is the statement. (ii) At $\hat\beta=0$, $g=X^\top y/n$ and (i) requires $|g_j|\le\lambda$ for
all $j$. (iii) With $X^\top X=nI$ the objective separates into $\tfrac12(\beta_j-x_{(j)}^\top y/n)^2+\lambda|\beta_j|$
plus a constant, one term for each $j$, and the minimiser of each is the soft threshold, by (i) in one dimension.
(iv) As a function of $\beta_j$ alone the objective is $\tfrac12c_j\beta_j^2-z_j\beta_j+\lambda|\beta_j|$ plus a
constant. Setting zero in its subdifferential gives $c_j\beta_j=z_j-\lambda$ if $\beta_j>0$,
$c_j\beta_j=z_j+\lambda$ if $\beta_j<0$, and $|z_j|\le\lambda$ if $\beta_j=0$, which together say
$\beta_j=S(z_j,\lambda)/c_j$.
\end{proof}

Part (iv) is the coordinate descent algorithm of \textcite{friedman2010}: sweep through the
coordinates, soft-threshold each in turn and update the residual. The listing is the whole
method. Its fixed number of sweeps keeps it short, where working code would also stop once the
coefficients stop changing, and it skips a column of zeros, whose coefficient stays at zero. The library version works on the Gram matrix $X^\top X/n$, which is faster when $n>p$, and a test
checks it against scikit-learn.
\Needspace*{24\baselineskip}
\lstinputlisting[language=Python]{code/ch12/snippets/lasso_cd.py}

\subsection{A simulation}

The script \texttt{regularise} draws \bvReps{} training samples\ix{Training set} of $n=\bvN$ observations on $p=\bvP$
predictors. The predictors are normal with correlation $\bvRho^{|j-k|}$ between columns $j$ and $k$, five
coefficients are nonzero, and the noise\ix{Noise} has standard deviation \bvSig, so that the signal variance is
\bvSnr{} of the noise variance. Each sample is fitted by least squares\ix{Least squares} and along a grid of penalties for ridge\ix{Ridge regression}
and the lasso\ix{Lasso}, and every fit predicts the same \bvNte{} test points. Averaging over the samples splits the
excess test error, the part above $\sigma^2$ in \cref{eq12:bv}, into squared bias\ix{Bias} and variance\ix{Bias-variance trade-off}
(\cref{fig12:bv}).

\begin{figure}[!tb]
  \centering
  \includegraphics[width=0.94\textwidth]{ch12/fig01_bias_variance}
  \caption{\kind{simulation} Excess test error of ridge (left) and the lasso (right) against the penalty, split into squared
  bias and variance, averaged over \bvReps{} simulated training samples of \bvN{} observations on \bvP{}
  correlated predictors. The horizontal line is least squares.}
  \label{fig12:bv}
\end{figure}

Least squares has an excess error of \bvOlsErr, of which \bvOlsVar{} is variance and \bvOlsBias{} the
simulation's estimate of a bias that is zero. Ridge at the penalty that is best in hindsight, $\lambda=\bvRidgeLam$ (an oracle choice that
cross-validation\ix{Cross-validation} can only approximate), cuts the
excess error to \bvRidgeErr, \bvRidgeCut\% less, with squared bias \bvRidgeBias{} and variance
\bvRidgeVar; its effective degrees of freedom are about \bvRidgeDf{} of the \bvP. Even the smallest
penalty on the grid, $\lambda=\bvRidgeSmallLam$, lowers the error a little, to \bvRidgeSmall, consistent with
\cref{prop12:ridge}(iii), which guarantees an improvement in the population risk for small enough
penalties, though not in every simulated sample. The lasso does better still here, \bvLassoErr{} at
$\lambda=\bvLassoLam$, because the true model is sparse: it keeps \bvLassoNz{} predictors on average,
including most of the five that matter. When many small coefficients all matter, ridge tends to win; when
a few large ones do, the lasso tends to. \Cref{fig12:paths} shows the coefficients of one sample along each
path. Ridge moves every coefficient smoothly towards zero; the lasso drops them one at a time.

\begin{figure}[!tb]
  \centering
  \includegraphics[width=0.94\textwidth]{ch12/fig02_paths}
  \caption{\kind{simulation} Coefficient paths for one simulated sample. The five coefficients that are nonzero in the true
  model are drawn in red. The lasso sets coefficients to zero one at a time as $\lambda$ grows; ridge only
  shrinks them.}
  \label{fig12:paths}
\end{figure}

How closely does cross-validation find the oracle? \Cref{fig12:cvoracle} answers with a separate
simulation of the same model, \cvReps{} training samples on random stream 13, in each of which the
penalty is chosen by \cvK-fold cross-validation. On average the cross-validation curve follows the excess
error of the fit near its minimum and lies above it for small penalties, where a fit on nine tenths of
the sample has more variance than the fit on all of it. The chosen ridge penalty has a median of
\cvRidgeMed, and its 10th and 90th percentiles are \cvRidgeLo{} and \cvRidgeHi, around the oracle
value of \cvRidgeOracle{} in this simulation; for the lasso the figures are \cvLassoMed, \cvLassoLo{} and
\cvLassoHi{} around \cvLassoOracle. The penalty chosen leaves an average excess error of \cvRidgeErrCv{}
for ridge and \cvLassoErrCv{} for the lasso, \cvRidgeExtra\% and \cvLassoExtra\% above the
\cvRidgeErrOracle{} and \cvLassoErrOracle{} of the oracle penalty in the same samples: cross-validation
gives away a little of the gain, and keeps most of it.

\begin{figure}[!tb]
  \centering
  \includegraphics[width=0.94\textwidth]{ch12/fig15_cv_oracle}
  \caption{\kind{simulation} Random stream 13: \cvK-fold cross-validation against the excess error of the
  fit for ridge (left) and the lasso (right), averaged over \cvReps{} training samples of the model of
  \cref{fig12:bv}. The cross-validation estimate is the out-of-fold mean squared error less the noise
  variance. The dotted line is the penalty with the smallest average excess error, and the shaded band
  runs from the 10th to the 90th percentile of the penalties that cross-validation chose.}
  \label{fig12:cvoracle}
\end{figure}

\FloatBarrier
% =====================================================================================
\section{How should a model be validated when time runs one way?}
\label{sec12:cv}

A penalty, a tree depth or a number of boosting rounds is chosen by estimating the test error of each
candidate. $K$-fold cross-validation splits the sample into $K$ folds, fits on $K-1$ of them and
predicts the one left out, in turn, and pools the errors. Its usual justification is that the left-out fold
behaves like new data, independent\ix{Independence} of the fitting sample. Financial data break that in two ways. Observations close in time
are dependent\ix{Dependence}, and a label often spans several periods: a twelve month return starting in January
shares\ix{Stock} eleven months with the one starting in February.

\begin{proposition}[Overlapping labels]\label{prop12:overlap}
Let $r_1,r_2,\dots$ be independent with common variance $s^2$, and let $y_t=r_{t+1}+\dots+r_{t+h}$.
\begin{enumerate}[label=(\roman*)]
\item $\operatorname{Corr}(y_t,y_{t+k})=(h-k)/h$ for $0\le k<h$, and zero for $k\ge h$.
\item For a test block of labels $t=a,\dots,b$, a training label $y_s$ shares no return with any test label if
and only if $s\le a-h$ or $s\ge b+h$. Removing the $h-1$ labels on each side of the block, $2(h-1)$ in all
(fewer for a block at either end of the sample), leaves the remaining training labels independent of
the test labels.
\end{enumerate}
\end{proposition}

\begin{proof}
(i) $y_t$ and $y_{t+k}$ share the returns $r_{t+k+1},\dots,r_{t+h}$, which are $h-k$ in number when $k<h$, so
$\operatorname{Cov}(y_t,y_{t+k})=(h-k)s^2$, while each has variance $hs^2$. (ii) $y_s$ uses $r_{s+1},\dots,r_{s+h}$
and the test labels together use $r_{a+1},\dots,r_{b+h}$. The two sets are disjoint exactly when $s+h\le a$ or
$s+1\ge b+h+1$. Labels built from disjoint sets of independent returns are independent.
\end{proof}

Removing the overlapping labels is called purging\ix{Cross-validation!purged}\index{Cross-validation!purged}, and dropping a further stretch after each test block,
to allow for dependence in the predictors, is called an embargo\index{Cross-validation!purged} \parencite{lopezdeprado2018}. The
simulated correlations of twelve month labels one, six and eleven months apart are \kfCorrOne, \kfCorrSix{}
and \kfCorrEleven, against \kfCorrThOne, \kfCorrThSix{} and \kfCorrThEleven{} from part (i).
\Cref{fig12:folds} shows the scheme.

\begin{figure}[!tb]
  \centering
  \includegraphics[width=0.9\textwidth]{ch12/fig03_purged_folds}
  \caption{The first five of \kfK{} purged and embargoed folds for \kfN{} twelve month labels. Each row is
  one fold: the test block, the labels purged on either side because their returns overlap the test labels'
  returns, and an embargo of \kfEmb{} labels after the purged labels that follow the block.}
  \label{fig12:folds}
\end{figure}

The listing below generates the folds. For a test block of labels $a,\dots,b$ and an embargo of $e$
labels it keeps a training label $s$ when $s\le a-h$ or $s\ge b+h+e$, the boundaries of part (ii) with
the embargo counted from the end of the purge, which is the convention of the chapter's library and of
\cref{fig12:folds}; an assertion checks that no kept label shares a return with a test label. For
\kfN{} labels each test block holds 46 or 47 labels and each training set 387 to 410, and the last line
checks the count asked for in \cref{exr12:four}.
\Needspace*{26\baselineskip}
\lstinputlisting[language=Python]{code/ch12/snippets/purged_kfold.py}

\subsection{What shuffling does}

The script \texttt{cv\_time} builds a world with nothing to find. Monthly returns are independent normal
with mean \kfMu\% and volatility \kfVol\%, and a predictor follows an autoregression with coefficient
\kfPhi, independent of the returns, like a valuation ratio. The label is the sum of the next \kfH{}
returns, and \kfT{} months give \kfN{} labels. A regression of the label on the predictor is scored by
three \kfK-fold schemes: folds drawn at random, contiguous blocks\ix{Block bootstrap}, and contiguous blocks purged\ix{Cross-validation!purged} and
embargoed. The fit from the whole sample is also scored on a fresh sample from the same model, which is
what the cross-validation is meant to estimate. Each scheme reports the pooled out-of-sample $R^2$, one
minus the ratio of the squared errors to those of the training mean.

\Cref{fig12:dp} sets the simulated predictor beside the kind of series it imitates. The log
dividend-price ratio of \cref{sec12:premium}, monthly from \dpFrom{} to \dpTo, has a first-order
autocorrelation of \dpRho, and its median monthly change is \dpMoveMed\% of its standard deviation, so
it drifts far from its mean and stays there for decades. One simulated path of the autoregression with
coefficient \kfPhi, scaled to the same mean and standard deviation, is less persistent\ix{Persistence}: its coefficient
gives a half-life of about \dpSimHalf{} years against about \dpHalf{} for the data, so it crosses its mean
several times a decade. The experiment below therefore understates the problem: a predictor as
persistent as the real ratio leaves even fewer independent observations in a sample.

\begin{figure}[!tb]
  \centering
  \includegraphics[width=0.94\textwidth]{ch12/fig19_dp}
  \caption{\kind{historical} Left: the log dividend-price ratio of the US market, monthly, \dpFrom{} to \dpTo, from the
  data of \textcite{welch2008}, updated by \textcite{goyalch12} and pulled by \texttt{scripts/pull\_ch12.py}. Right: simulation, random
  stream 51: one path of the autoregression with coefficient \kfPhi{} used by \texttt{cv\_time}, scaled
  to the mean and standard deviation of the data.}
  \label{fig12:dp}
\end{figure}

Over \kfReps{} runs the shuffled folds give a mean $R^2$ of \kfShufMean\%, positive in \kfShufPos\% of
runs and above 5\% in \kfShufHigh\%. The purged folds give \kfPurgeMean\%, and the fresh sample
\kfFreshMean\% (\cref{fig12:leak}). With random folds, almost every test label has neighbours in the
training set with nearly the same predictor value and a label sharing most of its returns; the regression
fits the accidental co-movement in the sample and is then scored on points that share it. Contiguous
blocks without purging, at \kfBlockMean\%, remove most of the effect, since only the labels at the
edges of each block overlap, and purging removes the rest.

\begin{figure}[!tb]
  \centering
  \includegraphics[width=0.86\textwidth]{ch12/fig04_cv_leak}
  \caption{\kind{simulation} Out-of-sample $R^2$ of a predictor with no power over twelve month returns, in \kfReps{}
  simulated samples: \kfK-fold cross-validation with shuffled folds, with purged and embargoed blocks, and
  the full-sample fit scored on a fresh sample.}
  \label{fig12:leak}
\end{figure}

\subsection{Leaks from outside the folds}

Any decision made with the whole sample leaks into every fold. A second experiment uses \kfSelT{} monthly
labels, which do not overlap, and \kfSelN{} unrelated predictors. Choosing the predictor most correlated with the
returns over the whole sample and then cross-validating it gives a mean $R^2$ of \kfSelOut\%, positive in
\kfSelOutPos\% of \kfSelReps{} runs. Repeating the choice inside each training fold, so that the test fold
plays no part in it, gives \kfSelIn\%, positive in \kfSelInPos\%. The same kind of leak, which this chapter does not simulate,
comes from standardising with full-sample means, filling gaps with later data, choosing the sample period after looking at results, and
using an index's current members to build its past. Cross-validation estimates the error of the whole
procedure only when the whole procedure is repeated inside each fold.

\FloatBarrier
% =====================================================================================
\section{What do trees, forests and boosting add?}
\label{sec12:trees}

A regression tree\ix{Decision tree} \parencite{breiman1984} splits the predictor space into rectangles and predicts the
mean response in each. It is grown greedily: at each node, every predictor $j$ and every threshold $t$
between adjacent values is tried, and the split that most reduces the sum of squared errors is kept,
until a depth or a minimum leaf size is reached.

\begin{proposition}[The best split]\label{prop12:split}
Split $n$ responses with sum $Y$ into a left group of $n_L$ with sum $Y_L$ and a right group with $n_R$
and $Y_R$. The fall in the sum of squared deviations from the group means is
$Y_L^2/n_L+Y_R^2/n_R-Y^2/n$.
\end{proposition}

\begin{proof}
For any group, $\sum(y_i-\bar y)^2=\sum y_i^2-(\sum y_i)^2/m$ with $m$ its size. The $\sum y_i^2$ terms are the
same before and after the split and cancel.
\end{proof}

So once the responses are sorted by a predictor, cumulative sums give every candidate split of that
predictor at once, and growing a tree costs little more than sorting. A tree captures interactions and
thresholds that a linear model misses, but a deep tree has high variance: a small change in the sample
moves an early split and everything below it.

The listing below is that inner loop. It sorts one predictor, forms the improvement of
\cref{prop12:split} at every cut from cumulative sums, skips cuts between tied values, and compares the
answer with a brute force search over all cuts on simulated data with a step at 0.45.
\Needspace*{26\baselineskip}
\lstinputlisting[language=Python]{code/ch12/snippets/best_split.py}

\subsection{Averaging and boosting}

Bagging \parencite{breiman1996} fits trees\ix{Decision tree} to bootstrap samples and averages them. A random forest\ix{Random forest}\index{Random forest}
\parencite{breiman2001} also lets each split consider only a random subset of $m$ of the $p$ predictors, which
makes the trees less alike. The next result says why that helps.

\begin{proposition}[Averaging correlated predictions]\label{prop12:avg}
If $\hat f_1,\dots,\hat f_B$ have common variance $s^2$ and common pairwise correlation $\rho$, their average has
variance $\rho s^2+(1-\rho)s^2/B$.
\end{proposition}

\begin{proof}
$\operatorname{Var}\bigl(\tfrac1B\sum_b\hat f_b\bigr)=\tfrac1{B^2}\bigl(Bs^2+B(B-1)\rho s^2\bigr)=\rho s^2+(1-\rho)s^2/B$.
\end{proof}

More trees remove the second term and leave the first, so the variance of a forest is limited by how
correlated its trees are, and the averaging does nothing for bias\ix{Bias}. Boosting\ix{Gradient boosting} attacks the bias instead.
For squared loss $\tfrac12(y-F)^2$, the negative gradient with respect to the current fit $F(x_i)$ is the
residual\ix{Residual} $y_i-F(x_i)$, and gradient boosting\index{Gradient boosting} \parencite{friedman2001} fits a small tree to the residuals,
adds a fraction $\nu$ of it to $F$, and repeats. Each step is a step of gradient descent\ix{Gradient descent} in the space of
functions, taken along the direction that a small tree can represent.

\begin{proposition}[Boosting a linear smoother]\label{prop12:boost}
Suppose every round fits the residuals with the same linear smoother, fitted values $Ay$ for responses $y$,
starting from $F_0=0$. After $M$ rounds with step $\nu$, $F_M=\bigl(I-(I-\nu A)^M\bigr)y$. If $A$ is
symmetric with eigenvalues\ix{Eigenvalue} $s_i\in[0,1]$ and $0<\nu\le1$, the fit multiplies the $i$th eigencomponent of $y$ by
$1-(1-\nu s_i)^M$, which rises to one as $M$ grows whenever $s_i>0$.
\end{proposition}

\begin{proof}
$F_{m+1}=F_m+\nu A(y-F_m)$, so $y-F_{m+1}=(I-\nu A)(y-F_m)$ and $y-F_M=(I-\nu A)^My$. The eigenvalue statement
follows from diagonalising $A$; $0\le1-\nu s_i\le1$.
\end{proof}

\Textcite{buhlmann2003} use this form to analyse boosting. Each extra round lowers the bias of every
component the smoother can see and raises the variance, so the number of rounds plays the part of the
penalty in ridge\ix{Ridge regression}, and a small $\nu$ makes the path finer. With trees the smoother changes from round to round
and the formula is only a guide, but the trade-off is the same.

\Cref{fig12:shrink} draws the trade-off for one eigencomponent. With $\nu=0.1$ boosting fits half of a
component with $s=0.9$ after \skHalfBig{} whole rounds and half of one with $s=0.05$ after \skHalfSmall, so stopping
early leaves the weak directions mostly unfitted, as a large ridge penalty does to directions with small
$d_j$; with $\nu=1$ every curve rises faster and the choice of when to stop is coarser.

\begin{figure}[!tb]
  \centering
  \includegraphics[width=0.94\textwidth]{ch12/fig16_shrink}
  \caption{\kind{exact} Left: the share $1-(1-\nu s)^M$ of an eigencomponent with eigenvalue $s$ fitted after $M$
  rounds of boosting a linear smoother at rate $\nu$. Right: the ridge factor $d^2/(d^2+\lambda)$ for
  three singular values. Both are exact formulas; no data or simulation is involved.}
  \label{fig12:shrink}
\end{figure}

\subsection{A benchmark}

The script \texttt{trees} uses the test function of \textcite{friedman1991},
$f(x)=10\sin(\pi x_1x_2)+20(x_3-\tfrac12)^2+10x_4+5x_5$, with \tbP{} predictors uniform on $[0,1]$, of which
five do not enter, noise\ix{Noise} of variance\ix{Variance} one, \tbN{} training points and \tbNte{} test points. The variance of
$f$ is \tbVarFPop, and \tbVarF{} over these test points. Errors below are mean squared errors against $f$ itself, so the
noise is excluded.

Linear regression has an error of \tbOls, since it cannot represent the sine or the square. A single tree
with its depth chosen by five-fold cross-validation, \tbDepth{} levels and \tbTreeLeaves{} leaves, does worse at
\tbTree; it can represent anything, but at this sample size the price is variance. Bagging\ix{Bagging} \tbB{} trees\ix{Decision tree} cuts
the error to \tbBag{} and a forest with \tbMtry{} of the \tbP{} predictors per split to \tbRf. The forest
is behind bagging here because half the predictors are noise, so about one split in twelve sees none of
the five useful predictors, and many see only the weaker ones. Boosting\ix{Gradient boosting} trees of depth three with $\nu=\tbNu$ has its lowest test error,
\tbBoostBest, after \tbBoostAt{} rounds, a minimum read off the test curve and so slightly flattering,
and is still at \tbBoost{} after \tbM. With $\nu=1$ it reaches only \tbBoostOne, after \tbBoostOneAt{} rounds,
and then gets worse (\cref{fig12:ens}). \Cref{fig12:treeboost} shows both effects in one dimension: a
deeper tree follows the noise, and more boosting rounds follow the curve.

To check \cref{prop12:avg}, \tbVarR{} training samples were drawn and a forest of \tbVarB{} trees grown on
each. At 200 fixed points, a single tree has variance \tbSigTwo{} across samples and randomisations on average,
and two trees from the same sample have correlation \tbRho, their average covariance over the average
variance. The proposition then gives a variance for the forest of \tbVarFormula, and the simulated
variance is \tbVarAvg. That is a small part of the forest's error of \tbRf, which suggests that most of
the rest is squared bias: the forest is the average of many trees, each of which is too coarse in the
directions where $f$ curves.

\begin{figure}[!tb]
  \centering
  \includegraphics[width=0.9\textwidth]{ch12/fig06_ensembles}
  \caption{\kind{simulation} Test error against the true function, by number of trees, for bagging, a random forest and
  boosting with two step sizes, on Friedman's benchmark with \tbN{} training points. The dotted lines are a
  single tree of cross-validated depth and linear regression.}
  \label{fig12:ens}
\end{figure}

\begin{figure}[!tb]
  \centering
  \includegraphics[width=0.94\textwidth]{ch12/fig05_tree_boost}
  \caption{\kind{simulation} One predictor and 200 noisy points. Left: trees of depth 2 and 5. Right: boosting trees of depth
  one with $\nu=\tbNu$, after 10 and 200 rounds.}
  \label{fig12:treeboost}
\end{figure}

\FloatBarrier
% =====================================================================================
\section{How is a neural network trained?}
\label{sec12:nn}

A feed-forward network with $L$ layers maps an input $a_0=x$ through
\begin{equation}\label{eq12:mlp}
  z_\ell=W_\ell^\top a_{\ell-1}+b_\ell,\qquad a_\ell=\phi(z_\ell)\quad(\ell<L),\qquad a_L=z_L ,
\end{equation}
where $\phi$ acts on each component; here $\phi=\tanh$ and the last layer is linear. With enough hidden
units such a network can approximate any continuous function on a compact set as closely as wanted
\parencite{cybenko1989,hornik1991}. That result says nothing about how many units are needed or how to find
the weights. The weights are found by minimising a loss $J$, here $\tfrac12(a_L-y)^2$ averaged over the
sample, with gradient steps, and the gradients come from the chain rule applied backwards through the layers.

\begin{proposition}[Backpropagation]\label{prop12:backprop}
For one observation, with $J=\tfrac12(a_L-y)^2$, let $\delta_\ell=\partial J/\partial z_\ell$. Then $\delta_L=a_L-y$,
\begin{equation}\label{eq12:backprop}
  \delta_{\ell-1}=\bigl(W_\ell\,\delta_\ell\bigr)\odot\phi'(z_{\ell-1}),\qquad
  \frac{\partial J}{\partial W_\ell}=a_{\ell-1}\,\delta_\ell^\top,\qquad \frac{\partial J}{\partial b_\ell}=\delta_\ell ,
\end{equation}
where $\odot$ multiplies component by component, and $\partial J/\partial x=W_1\delta_1$. For $\phi=\tanh$,
$\phi'(z_{\ell-1})=1-a_{\ell-1}^2$. For the loss averaged over a sample, each $\delta$ is divided by the
sample size, as the listing does.
\end{proposition}

\begin{proof}
$J$ depends on $W_\ell$, $b_\ell$ and $a_{\ell-1}$ only through $z_\ell$. Its components are
$z_{\ell,k}=\sum_iW_{\ell,ik}a_{\ell-1,i}+b_{\ell,k}$, so
\[
  \frac{\partial J}{\partial W_{\ell,ik}}=\delta_{\ell,k}a_{\ell-1,i},\qquad
  \frac{\partial J}{\partial b_{\ell,k}}=\delta_{\ell,k},\qquad
  \frac{\partial J}{\partial a_{\ell-1,i}}=\sum_kW_{\ell,ik}\delta_{\ell,k},
\]
the last being the $i$th component of $W_\ell\delta_\ell$. Since $a_{\ell-1,i}=\phi(z_{\ell-1,i})$, multiplying by $\phi'(z_{\ell-1,i})$
gives $\delta_{\ell-1,i}$. At the top, $a_L=z_L$ and $\partial J/\partial a_L=a_L-y$. The input case is the same
computation with $a_0=x$ and no activation, and $\tanh'=1-\tanh^2$.
\end{proof}

One backward pass therefore costs a small constant multiple of one forward pass, however many weights
there are,
which is why networks with many parameters can be trained at all \parencite{rumelhart1986}. The library
implements \cref{eq12:backprop} in a dozen lines, and a test compares every type of gradient with central
finite differences\ix{Finite differences}. Training uses mini-batches and the Adam update of \textcite{kingma2015}, which divides a
running mean of each gradient component by a running estimate of its root mean square.

The listing below writes \cref{eq12:backprop} out for one hidden layer of $\tanh$ units and the squared
error. On the network of \cref{exr12:five} it prints the fitted value $\hat y=0.9242$ and the gradients 0.4271 for $w_2$
and 1.4537 for $w_1$, which fixes the conventions, and on a network with two inputs and eight hidden
units it compares each of the 33 gradient components with a central finite difference.
\Needspace*{30\baselineskip}
\lstinputlisting[language=Python]{code/ch12/snippets/backprop.py}

\subsection{Learning the Black-Scholes price}

Can a network learn a pricing function from examples? \Textcite{hutchinson1994} asked this with simulated
and traded options. At a zero rate the call price\ix{Call price} per unit of strike\ix{Strike} depends only on the moneyness\ix{Moneyness} $m=S/K$
and the total volatility $v=\sigma\sqrt T$, through the formula of Chapter~5. The script \texttt{neural}
draws \nnNtr{} training pairs uniformly on $m\in[\nnMlo,\nnMhi]$ and $v\in[\nnVlo,\nnVhi]$, scales both
inputs to $[-1,1]$, and trains a network with two hidden layers of \nnHid{} units, \nnParams{} parameters in
all, for \nnEpochs{} passes through the data in batches of \nnBatch, \nnSteps{} steps, with a step size
falling from \nnLr{} to \nnLrEnd.

On a grid of test points the root mean square\ix{Mean squared error} error is \nnRmse{} basis points of the strike, and the
largest error is \nnMax, found where $v$ is small and the price has a sharp bend at the money
(\cref{fig12:nnerr}). A polynomial of degree \nnPolyDeg{} in the two scaled inputs, \nnPolyTerms{} terms fitted
by least squares on the same points, has errors of \nnPolyRmse{} and \nnPolyMax{}. The network's delta,
$\partial c/\partial m$, is available from the same backward pass by \cref{prop12:backprop}, taking the input
gradient and dividing by the scaling. Its root mean square error against $\Phi(d_1)$ is \nnDeltaRmse, and its
largest error, \nnDeltaMax, is at moneyness \nnDeltaMaxM{} and total volatility \nnDeltaMaxV, the corner of the training range where both are
highest. \Cref{fig12:nntrain} shows the same failure on the line $\sigma\sqrt T=0.1$, where the
network's delta falls below $\Phi(d_1)$ as $m$ approaches \nnMhi. The delta was not part of the training loss; here a small price error came with a
small delta error over most of the range, but nothing guarantees it, and the edges are where it fails first.

\begin{figure}[!tb]
  \centering
  \includegraphics[width=0.96\textwidth]{ch12/fig07_nn_error}
  \caption{\kind{simulation} Error of the fitted call price per unit of strike, in basis points of the strike, over moneyness
  and total volatility: the network with \nnParams{} parameters (left) and a polynomial of degree
  \nnPolyDeg{} with \nnPolyTerms{} terms (right).}
  \label{fig12:nnerr}
\end{figure}

\begin{figure}[!tb]
  \centering
  \includegraphics[width=0.96\textwidth]{ch12/fig08_nn_training}
  \caption{\kind{simulation} Left: mean squared error of the network on the training and validation points during training.
  Right: the network's delta, from backpropagation\index{Backpropagation} to the inputs, at two total volatilities, with the
  Black-Scholes delta.}
  \label{fig12:nntrain}
\end{figure}

Two cautions apply. The fitted price is not guaranteed to respect no-arbitrage: its second difference in
moneyness is negative at \nnConvex\% of the grid points, violations of the convexity\ix{Convexity} a call price
must have, concentrated\ix{Concentration} at the in-the-money edge, and a network used for pricing needs such constraints built in or checked. And the formula was
known, so the exercise measures the method, not the market. Its use in practice is to replace a slow pricer,
a Monte Carlo or a finite difference scheme, by a fast approximation trained on its output, and the errors
above are the kind of price that approximation costs.

\Cref{fig12:gamma} shows where the convexity fails. It compares a numerical second derivative of the
normalised call price $c=C/K$ with respect to moneyness $m=S/K$ against $K\Gamma$, where
$\Gamma=\partial^2C/\partial S^2$ is the Black-Scholes\ix{Black-Scholes formula} gamma\ix{Gamma}, along two lines of total volatility. At
$\sigma\sqrt T=0.1$ the network follows the peak of \gmPeakLow{} near the money with a root mean
square error of \gmErrLow; at $\sigma\sqrt T=0.4$, where the gamma never exceeds \gmPeakHigh, the
error is \gmErrHigh{} and the network's curve oscillates around it. On both lines the negative values,
\gmNegLow\% and \gmNegHigh\% of the points, sit at the deep-in-the-money edge of the training
region, where the network is held by data on one side only. On the low-volatility line the price there
is also nearly linear, with a gamma of \gmEdgeLow{} against the peak of \gmPeakLow, so nothing pins its
curvature; on the high-volatility line the gamma at the edge is still \gmEdgeHigh.
The band of \cref{sec12:hedge} uses the analytical Black-Scholes gamma; a band built from the
network's gamma would inherit this curvature error at the edge of the range.

\begin{figure}[!tb]
  \centering
  \includegraphics[width=0.94\textwidth]{ch12/fig17_nn_gamma}
  \caption{\kind{simulation} Numerical differentiation of the trained network: the second derivative of its normalised
  call price with respect to moneyness (solid), compared with $K$ times the Black-Scholes gamma (dotted),
  at two total volatilities. Negative second derivatives, marked, are failures of call-price convexity.}
  \label{fig12:gamma}
\end{figure}

\FloatBarrier
% =====================================================================================
\section{Can the equity premium be forecast out of sample?}
\label{sec12:premium}

A long literature reports variables that predict stock market returns in sample: valuation ratios,
interest rates\ix{Interest rate}, spreads, volatility and issuance. \Textcite{welch2008} asked whether any of them would
have helped an investor who had to estimate the forecasting regression with data available at the time,
and found that almost none did; \textcite{goyal2024} repeat the exercise with later data and more
variables. This section runs that test on the monthly data to December 2025 and adds the methods of the
previous sections.

\begin{definition}[Out-of-sample $R^2$]\label{def12:r2os}
For forecasts $\hat r_{t+1}$ and the historical mean\ix{Historical mean} $\bar r_t$, both formed with data through month $t$,
over an evaluation period $\mathcal T$,
\[
  R^2_{\mathrm{OS}}=1-\frac{\sum_{t\in\mathcal T}(r_{t+1}-\hat r_{t+1})^2}{\sum_{t\in\mathcal T}(r_{t+1}-\bar r_t)^2}.
\]
\end{definition}

The statistic of \textcite{campbell2008} is positive when the model beats the historical mean in mean
squared error\ix{Mean squared error}. It is the right comparison because the historical mean is what an investor with no model
would use. Testing whether it is significantly positive needs care, because the model nests the
benchmark\ix{Benchmark}: if the predictor is useless, the model still estimates a slope, and the noise in that estimate
makes the model's errors larger than the benchmark's on average, so a plain comparison of mean squared
errors is biased against the model.

\begin{proposition}[The Clark and West adjustment]\label{prop12:cw}
Let $e_{t+1}=r_{t+1}-\bar r_t$ and $d_{t+1}=e_{t+1}^2-\bigl[(r_{t+1}-\hat r_{t+1})^2-(\hat r_{t+1}-\bar r_t)^2\bigr]$. Then
\[
  e_{t+1}^2-(r_{t+1}-\hat r_{t+1})^2=d_{t+1}-(\hat r_{t+1}-\bar r_t)^2 ,\qquad d_{t+1}=2\,e_{t+1}(\hat r_{t+1}-\bar r_t).
\]
If the benchmark equals the true conditional mean, $\E_t r_{t+1}=\bar r_t$, then $\E d_{t+1}=0$, while the difference in squared errors has mean $-\E(\hat r_{t+1}-\bar r_t)^2\le0$.
\end{proposition}

\begin{proof}
Write $r-\hat r=e-(\hat r-\bar r)$ with time subscripts dropped. Then $(r-\hat r)^2=e^2-2e(\hat r-\bar r)+(\hat r-\bar r)^2$,
which gives both identities. Both forecasts are known at $t$, so $\E_t[e_{t+1}(\hat r_{t+1}-\bar r_t)]=
(\hat r_{t+1}-\bar r_t)\E_te_{t+1}=0$ under the hypothesis, and the law of iterated expectations gives $\E d_{t+1}=0$.
\end{proof}

The proposition is an idealised case: in practice $\bar r_t$ is an estimated historical mean, which the
test below allows for. \Textcite{clark2007} propose the $t$-statistic of the mean of $d_{t+1}$, compared with one-sided normal
critical values\ix{Critical value}, as an approximately correctly sized test of the hypothesis that the predictor is useless. The
approximation is derived for nested linear models; estimated means, expanding windows and serial
dependence affect its size in finite samples.

The R listing below runs the test for the Treasury bill\ix{Treasury bill} rate alone. It rebuilds the predictors from
the data file so that its sample is that of the Python code, forms an expanding-window least squares
forecast for each month from January 1965 using only the pairs known before it, and computes
$R^2_{\mathrm{OS}}$, the terms $d_{t+1}$, their $t$-statistic and the one-sided $p$ value. It prints
\eqOosN{} forecasts, $R^2_{\mathrm{OS}}=\eqRtwoTbl\%$, $t=\eqCwTbl$ and $p=\eqPTbl$, the figures of
\cref{tab12:premium}, and a test checks them.
\Needspace*{28\baselineskip}
\lstinputlisting[language=R]{code/ch12/snippets/clark_west.R}

\subsection{What a small $R^2$ is worth}

A monthly $R^2$ of half a percent sounds negligible. \Textcite{campbell2008} showed that it need not be.

\begin{proposition}[The value of a small $R^2$]\label{prop12:ct}
Let $r_{t+1}=\mu+x_t+\varepsilon_{t+1}$, where $x_t$ is known at $t$, $\E x_t=0$, $\operatorname{Var}x_t=R^2s^2$, and
$\varepsilon_{t+1}$ has mean zero and variance $(1-R^2)s^2$ given $x_t$. Let $\mathit{SR}_m=\mu/s$, with $\mu\ne0$ so that the ratio below is defined. A mean-variance investor
with risk aversion\ix{Risk aversion} $\gamma$ who knows $x_t$ holds $w_t=(\mu+x_t)/(\gamma(1-R^2)s^2)$ in the market; one who
does not holds $w=\mu/(\gamma s^2)$. The ratio of their average excess returns is
\begin{equation}\label{eq12:ct}
  \frac{\E[w_tr_{t+1}]}{\E[wr_{t+1}]}=\frac{\mathit{SR}_m^2+R^2}{\mathit{SR}_m^2(1-R^2)}=1+\frac{R^2}{\mathit{SR}_m^2}\,\frac{1+\mathit{SR}_m^2}{1-R^2}.
\end{equation}
\end{proposition}

\begin{proof}
Given $x_t$ the return has mean $\mu+x_t$ and variance $(1-R^2)s^2$, which gives $w_t$. Then
$\E[w_tr_{t+1}]=\E[(\mu+x_t)^2]/(\gamma(1-R^2)s^2)=(\mu^2+R^2s^2)/(\gamma(1-R^2)s^2)$, against
$\mu^2/(\gamma s^2)$ for the constant weight. Dividing and writing $\mu=\mathit{SR}_m\,s$ gives the first form; the second is
algebra.
\end{proof}

What matters is $R^2$ relative to $\mathit{SR}_m^2$, the squared monthly Sharpe ratio of the market, and that is small
too. From 1926 to 1964 the monthly Sharpe ratio of the log premium in the data is \eqS, \eqSAnn{} a year, so
by \cref{eq12:ct} an $R^2$ of 0.5\% would raise the timer's average excess return by \eqCtHalf\%. A test
checks \cref{eq12:ct} by simulation.

\Cref{fig12:ct} draws \cref{eq12:ct} for four monthly Sharpe ratios. Near $R^2=0$ the rise is approximately
linear, with exact initial slope $(1+\mathit{SR}_m^2)/\mathit{SR}_m^2$, so the same $R^2$ of 0.5\% that adds \eqCtHalf\% at the sample's
\eqS{} would add \ctLowHalf\% if the market's Sharpe ratio were 0.05 and much less if it were 0.15.

\begin{figure}[!tb]
  \centering
  \includegraphics[width=0.84\textwidth]{ch12/fig21_ct}
  \caption{\kind{exact} The proportional rise in a mean-variance timer's average excess return, \cref{eq12:ct},
  against the monthly $R^2$ of the predictor, for four monthly Sharpe ratios of the market. The mark
  is at $R^2=0.5\%$ for the sample's \eqS. The curves are exact; no data or
  simulation is involved.}
  \label{fig12:ct}
\end{figure}

\subsection{The data and the forecasts}

The fourteen predictors follow \textcite{welch2008}: the dividend\ix{Dividends}-price ratio\ix{Dividend-price ratio} and the dividend yield (the
latter on last month's price), the earnings-price ratio, the payout ratio, stock variance, the book to
market ratio, net equity expansion, the Treasury bill\ix{Treasury bill} rate, the long-term yield and return, the term spread,
the default\ix{Default} yield spread\ix{Yield spread}, the default return spread and inflation, lagged a month further because it is
published late. The target is the monthly log return of the S\&P 500 with dividends less the log Treasury
bill rate. Two of the predictors are exact combinations of others (the payout ratio is the
dividend-price ratio less the earnings-price ratio, and the term spread the long yield less the bill rate),
so the slopes of the regression on all fourteen span at most twelve independent directions, and
$X^\top X$ is singular. The code checks that the design with its constant has rank thirteen in every
window and fits it by \texttt{numpy.linalg.lstsq}, which returns the minimum-norm solution; because the
two combinations hold exactly out of sample too, any other solution gives the same forecasts. The first observation with
every predictor is in \eqFirst. Forecasts start in January \eqOosFrom{} and run to December \eqOosTo,
\eqOosN{} months, each made with an expanding window of the data available at the time.

The models are the historical mean\ix{Historical mean}; each predictor alone; all fourteen by least squares, with and without
flooring the forecast at zero as \textcite{campbell2008} suggest; the mean of the fourteen single-predictor
forecasts \parencite{rapach2010}; ridge and the lasso on standardised predictors, with the penalty chosen
each January to minimise the squared error, over the last \eqValid{} months of the window, of
forecasts from a fit on the months before them; a random forest\ix{Random forest} of \eqForestB{} trees\ix{Decision tree}, each split choosing among
\eqForestMtry{} predictors, with at most eight levels and at least \eqLeaf{} months in each leaf; and boosting with trees of
depth two, $\nu=\eqBoostNu$ and \eqBoostM{} rounds. The trees are refitted each January. The investor has
risk aversion \eqGamma, estimates the variance from the previous \eqVarWin{} months, and holds between 0 and
\eqWhi{} times wealth in the market; the certainty equivalent gain is the annualised mean less $\tfrac{\gamma}{2}$ times the annualised
variance of the model's portfolio, minus the same quantity for the portfolio of the historical mean.
Standardisation, lags and the choice of penalty use only data available at each forecast date, and
the economic results exclude transaction costs\ix{Transaction costs}.

\begin{table}[!tb]
\centering
\caption{\kind{historical} Monthly forecasts of the US equity premium, January \eqOosFrom{} to December \eqOosTo, against the
historical mean. Three of the fourteen single predictors are shown; $R^2_{\mathrm{OS}}$ for all fourteen is in \cref{fig12:r2os}.
$R^2_{\mathrm{OS}}$ over the whole period and before and after 2000, in percent; the Clark and
West statistic and its one-sided $p$-value, a formal test for the nested linear forecasts and a
descriptive diagnostic for the forest and boosting; the certainty equivalent gain for a mean-variance investor, in
percent a year; and the annualised Sharpe ratio of the investor's portfolio, which is \eqSrHm{} for the
historical mean.}
\label{tab12:premium}
\begin{tabular}{lrrrrrrr}
\toprule
 & \multicolumn{3}{c}{$R^2_{\mathrm{OS}}$ (\%)} & & & & \\
\cmidrule(lr){2-4}
model & all & to 1999 & from 2000 & CW & $p$ & CE gain & Sharpe\\
\midrule
Treasury bill rate      & \eqRtwoTbl & \eqRtwoATbl & \eqRtwoBTbl & \eqCwTbl & \eqPTbl & \eqCeTbl & \eqSrTbl\\
dividend-price ratio    & \eqRtwoDp & \eqRtwoADp & \eqRtwoBDp & \eqCwDp & \eqPDp & \eqCeDp & \eqSrDp\\
net equity expansion    & \eqRtwoNtis & \eqRtwoANtis & \eqRtwoBNtis & \eqCwNtis & \eqPNtis & \eqCeNtis & \eqSrNtis\\
all, least squares      & \eqRtwoOls & \eqRtwoAOls & \eqRtwoBOls & \eqCwOls & \eqPOls & \eqCeOls & \eqSrOls\\
all, floored at zero    & \eqRtwoOlsct & \eqRtwoAOlsct & \eqRtwoBOlsct & \eqCwOlsct & \eqPOlsct & \eqCeOlsct & \eqSrOlsct\\
mean of forecasts       & \eqRtwoComb & \eqRtwoAComb & \eqRtwoBComb & \eqCwComb & \eqPComb & \eqCeComb & \eqSrComb\\
ridge                   & \eqRtwoRidge & \eqRtwoARidge & \eqRtwoBRidge & \eqCwRidge & \eqPRidge & \eqCeRidge & \eqSrRidge\\
lasso                   & \eqRtwoLasso & \eqRtwoALasso & \eqRtwoBLasso & \eqCwLasso & \eqPLasso & \eqCeLasso & \eqSrLasso\\
random forest           & \eqRtwoForest & \eqRtwoAForest & \eqRtwoBForest & \eqCwForest & \eqPForest & \eqCeForest & \eqSrForest\\
boosting                & \eqRtwoBoost & \eqRtwoABoost & \eqRtwoBBoost & \eqCwBoost & \eqPBoost & \eqCeBoost & \eqSrBoost\\
\bottomrule
\end{tabular}
\end{table}

\subsection{Results}

\begin{figure}[!tb]
  \centering
  \includegraphics[width=0.86\textwidth]{ch12/fig10_r2os}
  \caption{\kind{historical} Out-of-sample $R^2$ against the historical mean, January \eqOosFrom{} to December \eqOosTo, for each
  predictor alone (above the dotted line) and for the combined models (below). Blue bars are positive; no
  model is significant at 5\% by the Clark and West test. Bars beyond $-4\%$ are cut and labelled.}
  \label{fig12:r2os}
\end{figure}

Only one of the fourteen single predictors beats the historical mean\ix{Historical mean} over the whole period, the
\eqBestUni, with $R^2_{\mathrm{OS}}=\eqBestUniRtwo\%$, and none is significant at 5\% by the Clark and West\ix{Clark-West test} test;
the worst, the \eqWorstUni, has \eqWorstUniRtwo\% (\cref{fig12:r2os}). Least squares\ix{Least squares} on all fourteen has
$R^2_{\mathrm{OS}}=\eqRtwoOls\%$, the variance of twelve estimated slopes at work, and flooring the forecast at
zero cuts the loss to \eqRtwoOlsct\% without removing it (\cref{tab12:premium}). The mean of the fourteen forecasts is the one combined
model that is positive over the whole period, \eqRtwoComb\%, with a Clark and West $p$-value of \eqPComb, because
averaging shrinks each forecast towards the historical mean, much as ridge\ix{Ridge regression} shrinks coefficients, and
diversifies the errors of the single models. Ridge itself is positive before 2000, \eqRtwoARidge\%, and negative after, \eqRtwoBRidge\%; the
lasso keeps \eqLassoNz{} predictors on average and does no better. The forest and boosting, which can fit
nonlinear and interacting effects, do worst of the regularised methods, at \eqRtwoForest\% and
\eqRtwoBoost\%: with a monthly signal this weak and between \eqWinFirst{} and \eqWinLast{} monthly observations in the
window, flexibility is mostly
variance, as in \cref{sec12:reg}. \Textcite{gu2020} report gains from trees and networks in the cross-section
of individual stocks, where thousands of observations a month give the methods more to work with. The
aggregate premium offers one observation a month.

\Cref{fig12:cumsse} shows where the results come from. The cumulative difference in squared errors
rises when the model is beating the mean and falls when it is not. The dividend-price ratio built up gains until late
1990 and lost them in the decade after. Ridge peaked at the same time, gave back most of its gain by
2010 and then lost far more in a few months of 2020. The forest was below zero for most of the period
from the mid-1970s, rising briefly above it around 1990, and lost ground in a few large steps. The mean
of the forecasts peaked early in 2009, and least squares on all fourteen lost ground in every decade but the \eqOlsGainDecade, when it
gained a little back, with its worst years in \eqOlsWorstYear{} and \eqOlsSecondYear.

\begin{figure}[!tb]
  \centering
  \includegraphics[width=0.9\textwidth]{ch12/fig09_cum_sse}
  \caption{\kind{historical} Cumulative squared forecast error of the historical mean less that of each model. The line rises
  in months where the model forecasts better.}
  \label{fig12:cumsse}
\end{figure}

The certainty equivalent column tells a slightly different story. Several models with negative
$R^2_{\mathrm{OS}}$ have positive gains, such as least squares at \eqCeOls\% a year and ridge at \eqCeRidge\%,
because the investor's weight is capped at 0 and \eqWhi: the worst forecasts are cut off, and the remaining
signal is in the sign more than the size. Least squares and its floored version have the same gain
because the cap at zero already floors the investor's weight. Forecast accuracy and portfolio value are related through
\cref{prop12:ct} only when the investor can act on the forecast without limits. The gains are also spread over
\eqOosN{} months with large sampling error, a point \cref{sec12:overfit} returns to.

\FloatBarrier
% =====================================================================================
\section{Learning a hedge}
\label{sec12:hedge}

\Cref{sec12:nn} fitted a network to prices. A different use of the same tools is to learn the hedge itself:
choose a rule that maps what the hedger knows to a position, and train it on simulated paths to minimise a
measure of the hedger's loss. \Textcite{buehler2019} call this deep hedging\ix{Deep hedging}\ix{Hedge} and use it where no formula
exists, with costs, many instruments and constraints. This section applies it where the answer is known,
a call on a lognormal stock, so that the learned hedge can be checked, and then adds costs, where it is not.

\begin{proposition}[Quadratic hedging is a regression]\label{prop12:qhedge}
Let $S_0,\dots,S_n$ be a martingale\ix{Martingale} under the probability measure of the simulation, with finite second moments\ix{Second moment}, and let the hedge ratio be $\Delta_k=\phi(S_k,t_k)^\top\theta$ for a
vector of features $\phi$ with finite second moments of $\phi_j(S_k,t_k)(S_{k+1}-S_k)$. Write
$G_j=\sum_{k<n}\phi_j(S_k,t_k)(S_{k+1}-S_k)$ for the gain of feature $j$, so the hedge gains $G^\top\theta$.
\begin{enumerate}[label=(\roman*)]
\item $\E G_j=0$ for every $j$, so $\E[G^\top\theta]=0$ for every $\theta$.
\item The pair $(p,\theta)$ that minimises $\E\bigl[(H-p-G^\top\theta)^2\bigr]$ for a claim $H$ with finite variance
is the population least squares\ix{Least squares} regression\ix{Linear regression} of $H$ on a constant and $G$. Its intercept is $p=\E H$, and $\theta$
minimises $\operatorname{Var}(H-G^\top\theta)$.
\item On a sample of paths, ordinary least squares of $H$ on a constant and $G$ gives the sample analogue,
with $\hat p=\bar H-\bar G^\top\hat\theta$.
\end{enumerate}
\end{proposition}

\begin{proof}
(i) $\phi_j(S_k,t_k)$ is known at $k$, so $\E[\phi_j(S_k,t_k)(S_{k+1}-S_k)]=\E\bigl[\phi_j(S_k,t_k)\,\E_k(S_{k+1}-S_k)\bigr]=0$;
sum over $k$. (ii) The objective is a least squares problem in $(p,\theta)$ with regressors $1$ and $G$. The
first-order condition for $p$ is $p=\E H-\E G^\top\theta=\E H$ by (i). Substituting, the objective is
$\operatorname{Var}(H-G^\top\theta)$. (iii) The same normal equations with sample moments.
\end{proof}

The paths here are simulated with zero drift and a zero interest rate, so the simulation's measure is
itself a risk-neutral\ix{Risk-neutral pricing} (martingale) measure and $\E H$ is the risk-neutral price; the regression then
learns both the price and the hedge. Under a physical measure with a drift, the same regression would
give an expectation under that measure, not a price. Under the risk-neutral measure the regression is the discrete form of the replication in Chapter~5, where the
discrete hedging\ix{Discrete hedging} error had mean zero and a variance proportional to the time step.

\subsection{Without costs}

The stock starts at \hdSZero, has volatility \hdSig\% and zero drift, and the call has strike \hdK{} and
three months to run, rebalanced daily\ix{Discrete hedging} over \hdSteps{} steps. The Black and Scholes price is \hdBs. Since
$d_1=z/\sigma+\sigma\sqrt\tau/2$ with $z=\ln(S/K)/\sqrt\tau$, the delta is a function of $z$ and $\sqrt\tau$ only,
and the features are chosen to let a linear rule reach it: a constant, \hdBumps{} logistic steps
$1/(1+e^{-(z-c)/w})$ with centres $c$ from $-\hdCentre$ to \hdCentre{} and width $w=\hdWidth$, and the same steps
multiplied by $\sqrt\tau$. That is \hdParams{} weights, shared by every date, the output layer of a network
whose hidden layer is fixed. The weights are fitted by \cref{prop12:qhedge} on \hdNtr{} paths and tested on
\hdNte{} fresh ones.

The learned price is \hdPriceLearn, against \hdBs, and the learned hedge ratio is within \hdDeltaGap{} of the
delta in root mean square over the test paths (\cref{fig12:hedge}). The standard deviation of the
hedger's loss is \hdStdBs{} for the delta and \hdStdLearn{} for the learned rule, against \hdUnhedged{} with no hedge, and it is
\hdStdLearnTrain{} on the training paths, so these diagnostics show little sign of overfitting. The
regression recovers the textbook answer from simulated paths with no knowledge of the model beyond the
choice of features, which used the form of $d_1$; the same code applies to paths from any simulator.

\begin{figure}[!tb]
  \centering
  \includegraphics[width=0.92\textwidth]{ch12/fig11_hedge_ratio}
  \caption{\kind{simulation} The learned hedge ratio without costs and the Black and Scholes delta after 20 and 55 of the
  \hdSteps{} trading days, with the learned no-trade band for costs of \hdKappa\%. The two hedge ratios are
  hard to tell apart.}
  \label{fig12:hedge}
\end{figure}

\subsection{With costs}

Now every trade costs \hdKappa\% of its value, including setting up the hedge and unwinding it at expiry.
\Cref{prop12:qhedge} no longer applies: the cost depends on the trades, and so on $\theta$, which makes the
loss non-quadratic in $\theta$, and a regression that ignores it chooses its weights as if trading
were free. Write $\mathcal L=H-G^\top\theta+C$ for the hedger's loss with costs $C$, and price by the
mean-variance\ix{Mean-variance optimisation} rule $\E\mathcal L+\tfrac\gamma2\operatorname{Var}\mathcal L$ with $\gamma=\hdGamma$. Under the delta the expected
cost\ix{Expected cost} is \hdCostBs, and the price is \hdPriceBs, against the Black and Scholes price of \hdBs; without
costs the same rule would give \hdPriceFree.
The weights learned without costs do no better, at \hdPriceZero.

Re-optimising the \hdParams{} weights for the mean-variance price, with L-BFGS on a smoothed
cost (the code checks the gradient against finite differences\ix{Finite differences}), lowers the price to \hdPriceCost{} after
\hdIters{} iterations: it trims the expected cost to \hdCostCost{} and the turnover\ix{Turnover} from \hdTurnBs{} to
\hdTurnCost{} shares per share hedged, and cuts the markup over the Black and Scholes price by \hdPriceCutCost\%.
The gain is small for a reason. A rule that depends only on the current price and date must trade whenever
the price moves, because its target moves; what saves costs is not trading when the position is close
enough to the target, and that needs to know the current position.

The policy of \textcite{whalley1997} has that memory. It holds the position while it lies within a band
around the delta and otherwise trades to the nearer edge. For small proportional costs, exponential
utility\ix{Utility} and a zero interest rate their asymptotic analysis gives the half-width
\begin{equation}\label{eq12:ww}
  \Bigl(\frac{3\kappa S\Gamma^2}{2\gamma}\Bigr)^{1/3},
\end{equation}
with $\Gamma$ the Black and Scholes gamma and $\kappa$ the proportional cost; at the start of this option it
is \hdWwAtm. The width is derived under those asymptotics and that utility, so the discrete
mean-variance experiment here uses it as a benchmark policy rather than an exact optimum. Learning a single scale for the band on the training paths gives \hdBandWw{} times
\cref{eq12:ww}, and on fresh paths the price falls to \hdPriceWw, a cut of \hdPriceCutWw\% in the markup, with
turnover \hdTurnWw{} and expected cost \hdCostWw. The loss is more variable, with standard deviation
\hdCStdWw{} against \hdCStdBs{} for the delta, and the mean-variance rule judges the saving in cost worth it
(\cref{fig12:hedgepnl}). The band at the asymptotic width, scale one, gives \hdPriceWwOne, and a band of
constant half-width \hdBandFlat{} gives \hdPriceFlat. Here a narrower band does a little better, but most of the gain comes from having a band at all.

\begin{figure}[!tb]
  \centering
  \includegraphics[width=0.82\textwidth]{ch12/fig12_hedge_pnl}
  \caption{\kind{simulation} The hedger's loss at expiry, costs included, on \hdNte{} fresh paths: the delta rebalanced daily
  and the no-trade band with the learned width. The band shifts the distribution to the left and widens it.}
  \label{fig12:hedgepnl}
\end{figure}

\Cref{fig12:hedgepath} follows one of the fresh paths. The band policy holds its position while the
delta moves inside the band and trades to the nearer edge when it leaves, so it trades on \hpTrades{} of
the \hpDays{} days and turns over \hpTurnBand{} shares against \hpTurnDelta{} for the daily delta.
\Cref{tab12:hedge} collects the six policies with costs. Every cell comes from the same run on the
fresh paths.

\begin{figure}[!tb]
  \centering
  \includegraphics[width=0.88\textwidth]{ch12/fig18_hedge_path}
  \caption{\kind{simulation} One fresh path of the hedging experiment with costs of \hdKappa\%. The shaded
  region is the learned no-trade band, \hdBandWw{} times \cref{eq12:ww} around the Black-Scholes delta;
  the step line is the band policy's position, and the marks are its trades.}
  \label{fig12:hedgepath}
\end{figure}

\begin{table}[!tb]
  \centering
  \small
  \caption{\kind{simulation} The six hedging policies with proportional costs of \hdKappa\%, on the same
  \hdNte{} fresh simulated paths: the mean-variance price, its markup over the Black-Scholes value of
  \hdBs, the expected cost, the turnover in shares and the standard deviation of the hedger's loss.}
  \label{tab12:hedge}
  \begin{tabular}{lrrrrr}
    \toprule
    policy & price & markup & cost & turnover & s.d.\ of loss\\
    \midrule
    Black-Scholes delta, daily & \hdPriceBs & \hdMarkupBs & \hdCostBs & \hdTurnBs & \hdCStdBs\\
    weights learned without costs & \hdPriceZero & \hdMarkupZero & \hdCostZero & \hdTurnZero & \hdCStdZero\\
    \hdParams{} weights learned with costs & \hdPriceCost & \hdMarkupCost & \hdCostCost & \hdTurnCost & \hdCStdCost\\
    band, scale 1 & \hdPriceWwOne & \hdMarkupWwOne & \hdCostWwOne & \hdTurnWwOne & \hdCStdWwOne\\
    band, learned scale \hdBandWw & \hdPriceWw & \hdMarkupWw & \hdCostWw & \hdTurnWw & \hdCStdWw\\
    constant band \hdBandFlat & \hdPriceFlat & \hdMarkupFlat & \hdCostFlat & \hdTurnFlat & \hdCStdFlat\\
    \bottomrule
  \end{tabular}
\end{table}

The lesson for learned hedges is about the state, more than the approximator: a network of any size fed\ix{Federal Reserve}
only the price and the date could not have found the band. When there are costs, the networks of
\textcite{buehler2019} take the previous position as an input.

\FloatBarrier
% =====================================================================================
\section{What a search does to a backtest}
\label{sec12:overfit}

Every method in this chapter has settings, and every setting tried and discarded is a trial. Chapter~8
gave the standard error of a Sharpe ratio and Chapter~7 the false discovery rate for many managers. This
section asks how much the best of many trials should be discounted for having been chosen.

\begin{proposition}[The best of $N$ trials]\label{prop12:emax}
Let $Z_1,\dots,Z_N$ be standard normal, independent or not. Then $\E\max_iZ_i\le\sqrt{2\ln N}$. If they are
independent,
\[
  \E\max_iZ_i=\int_{-\infty}^\infty xN\varphi(x)\Phi(x)^{N-1}\,dx ,
\]
and the ratio of this to $\sqrt{2\ln N}$ tends to one as $N$ grows.
\end{proposition}

\begin{proof}
For $t>0$, Jensen's inequality and $e^{t\max_iZ_i}\le\sum_ie^{tZ_i}$ give
$e^{t\E\max_iZ_i}\le\E e^{t\max_iZ_i}\le\sum_i\E e^{tZ_i}=Ne^{t^2/2}$, so $\E\max_iZ_i\le(\ln N)/t+t/2$ for every
$t>0$. The right side is smallest at $t=\sqrt{2\ln N}$, where it equals $\sqrt{2\ln N}$. With independence the
maximum has distribution function $\Phi(x)^N$ and density $N\varphi(x)\Phi(x)^{N-1}$, which gives the integral.
The limit of the ratio is a classical result of extreme value theory.
\end{proof}

\Textcite{bailey2014} use the approximation
$(1-\gamma_E)\Phi^{-1}(1-1/N)+\gamma_E\Phi^{-1}(1-1/(Ne))$, with $\gamma_E$ Euler's constant. For $N=10$, 100 and
1000 the exact values are \ovMaxTen, \ovMaxHundred{} and \ovMaxThousand, the approximation gives
\ovApproxTen, \ovApproxHundred{} and \ovApproxThousand, \ovSimReps{} simulated draws give \ovSimTen,
\ovSimHundred{} and \ovSimThousand, and the bound is \ovBoundTen, \ovBoundHundred{} and \ovBoundThousand{}
(\cref{fig12:emax}). A Sharpe ratio estimated over $Y$ years has a standard error of about $1/\sqrt Y$ in annual
units when the true ratio is near zero. So the best of 100 independent strategies with no skill, each
tested over \ovYears{} years, has an expected annual Sharpe ratio of \ovMaxSrHundred.

The Octave listing below recomputes the three columns in base Octave: the exact value by adaptive
quadrature of $x\,N\varphi(x)\Phi(x)^{N-1}$, the approximation through $\Phi^{-1}(p)=-\sqrt2\,
\mathrm{erfcinv}(2p)$, and the bound. It prints the values above and, for \cref{exr12:nine}, the
expected maximum\ix{Expected maximum} of 100 trials and the expected best annual Sharpe ratio over five years.
\Needspace*{17\baselineskip}
\lstinputlisting[language=Octave]{code/ch12/snippets/expected_max.m}

\begin{figure}[!tb]
  \centering
  \includegraphics[width=0.82\textwidth]{ch12/fig13_expected_max}
  \caption{\kind{exact} The expected maximum of $N$ independent standard normals, in units of the standard error of each
  trial: exact, the approximation of \textcite{bailey2014}, simulation, and the bound of \cref{prop12:emax}.}
  \label{fig12:emax}
\end{figure}

\subsection{A search on noise}

Simulate \ovT{} months of returns with a volatility of \ovVol\% a month and no predictability\ix{Return predictability} at all, and
search \ovRules{} trading rules\ix{Trading rule}: be long, or long and short, when the price is above its moving average over
2 to 51 months, and the two contrarian versions. Choose the rule with the best Sharpe ratio in the
first \ovHalf{} months after the longest average is available, and test it on the remaining \ovHalf. Over
\ovReps{} repetitions the chosen rule's annual Sharpe ratio averages \ovBestIs{} in sample and \ovBestOos{}
out of sample\ix{Out-of-sample}, with a standard error of \ovBestOosSe. The rules are correlated, with a median absolute
correlation of \ovCorr, and by \cref{prop12:emax} the in-sample maximum is about what \ovNeff{} independent trials
would give, a descriptive effective count and not a separately estimated number of tests.

\Textcite{bailey2017} turn the split into a statistic. Cut the sample into an even number $Q$ of blocks\ix{Block bootstrap}, and for each of the
$\binom{Q}{Q/2}$ ways of choosing half of them as the training set, pick the best rule in the training blocks
and find its relative rank $\omega\in(0,1)$ among all rules in the other half. The probability of backtest
overfitting\ix{Overfitting} is the share of splits in which the chosen rule ranks at or below the median out of sample, that is
with logit $\ln(\omega/(1-\omega))\le0$. With no skill anywhere the chosen rule is a random draw out of sample,
and the probability should be near one half. With $Q=\ovSRep$ it averages \ovPbo{} over \ovPboRuns{}
repetitions. \Cref{fig12:pbo} shows one run with $Q=\ovSFig$ and \ovSplits{} splits: the probability there is
\ovPboFig, and a regression of the chosen rule's out-of-sample Sharpe ratio on its in-sample ratio
has a slope of \ovSlope. The negative slope here is consistent with selection and with the way the procedure divides a
fixed sample: each rule's full-sample performance is fixed, so the blocks that flatter it in sample are
exactly the ones missing from its out-of-sample half. The procedure does not force a negative slope in
every case.

\begin{figure}[!tb]
  \centering
  \includegraphics[width=0.94\textwidth]{ch12/fig14_pbo}
  \caption{\kind{simulation} Combinatorially symmetric cross-validation of \ovRules{} rules on one simulated price with no
  predictability, with \ovSFig{} blocks and \ovSplits{} splits. Left: in-sample and out-of-sample annual Sharpe
  ratios of the chosen rule. Right: the logit of its relative rank out of sample; the share at or left of
  zero is the probability of backtest overfitting\index{Overfitting}.}
  \label{fig12:pbo}
\end{figure}

\subsection{Deflating a Sharpe ratio}

The probabilistic Sharpe ratio\ix{Sharpe ratio}\ix{Probabilistic Sharpe ratio}\index{Probabilistic Sharpe ratio} of \textcite{bailey2012} is
$\Phi\bigl((\hat\theta-\theta^*)/\hat\sigma(\hat\theta)\bigr)$, a normal-approximation score for the evidence
that the underlying ratio exceeds a benchmark $\theta^*$, and not a literal posterior probability, which
would need a Bayesian\ix{Bayes' rule} model; with $\hat\sigma(\hat\theta)$ the standard error of Chapter~8 that allows for skewness\ix{Skewness} and
kurtosis\ix{Kurtosis} \parencite{lo2002,opdyke2007}. The deflated Sharpe ratio\ix{Deflated Sharpe ratio} of \textcite{bailey2014} sets $\theta^*$ to the
expected maximum\ix{Expected maximum} of $N$ trials with no skill, the approximation above times the standard deviation of the
Sharpe ratios across the trials; for correlated trials chosen adaptively that adjustment is a heuristic.

Apply it to \cref{sec12:premium}. Each of the \ovEqModels{} models defines a timing portfolio, and its gain
over the portfolio of the historical mean is a monthly return series over the same \ovEqT{} months, with the same definition for every model. Least squares and its floored version hold the same portfolio,
since a negative forecast gives a zero weight either way, so there are \ovEqDistinct{} distinct series. The best of them
is the \ovEqBest, with a monthly Sharpe ratio of \ovEqSr{} on the difference, \ovEqSrAnn{} a year, skewness
\ovEqSkew{} and excess kurtosis \ovEqKurt. Taken alone the ratio has a standard error of \ovEqSe{} and a probabilistic
Sharpe ratio of \ovEqPsr{} against zero. The Sharpe ratios of the \ovEqModels{} models have a standard deviation
of \ovEqSrSd, so the benchmark is \ovEqBench, and the deflated Sharpe ratio is \ovEqDsr. Counting the two identical series once, $N=\ovEqDistinct$
with the spread over the distinct series, gives \ovEqDsrDistinct; the conclusion does not change. Once the search is counted, the
diagnostic falls well below 0.95. The models are correlated, so \ovEqModels{} overstates the number of independent trials, but the
count leaves out the searches that produced the list of predictors in the first place; the deflated
ratio is therefore a sensitivity\ix{Sensitivity analysis} measure and not an exact probability after selection.

\Cref{fig12:dsr} is that sensitivity. With the same Sharpe ratio, moments and cross-trial spread, the
deflated ratio is \dsTwo{} for two trials, falls below 0.95 from $N=\dsCross$, is \dsAtModels{} at the
\ovEqModels{} models and \dsHundred{} at 100. The verdict therefore rests less on the exact count than on
whether there were more than a handful of trials, and there were.

\begin{figure}[!hb]
  \centering
  \includegraphics[width=0.84\textwidth]{ch12/fig20_dsr}
  \caption{\kind{historical} The deflated Sharpe ratio of the best timing portfolio, the \ovEqBest, against the assumed
  number of independent trials $N$, with the observed Sharpe ratio, skewness, excess kurtosis and cross-trial
  standard deviation held at their sample values. The dashed line is 0.95; the dotted line marks the
  \ovEqModels{} models.}
  \label{fig12:dsr}
\end{figure}

\subsection{Many tests at once}

The single-predictor tests of \cref{sec12:premium} are fourteen tests of the same question. Chapter~7 used
the rule of \textcite{benjamini1995}, which controls the expected share of false discoveries. The rule of
\textcite{holm1979} controls the familywise error\ix{Family-wise error rate}, the probability of any false rejection, and rejects every hypothesis that the
Bonferroni test rejects, often more.

\begin{proposition}[Holm's procedure]\label{prop12:holm}
Sort the $m$ p-values\ix{p value} as $p_{(1)}\le\dots\le p_{(m)}$, find the smallest $k$ with $p_{(k)}>\alpha/(m-k+1)$, and
reject the hypotheses with the $k-1$ smallest p-values (all of them if there is no such $k$). If each p-value of a
true hypothesis is uniform or stochastically larger, the probability of rejecting any true hypothesis is
at most $\alpha$, whatever the dependence between the tests.
\end{proposition}

\begin{proof}
Let $I_0$ be the true hypotheses, $m_0=|I_0|\ge1$, and let $j$ be the rank of the smallest p-value among them.
At most $m-m_0$ hypotheses come before it, so $j\le m-m_0+1$ and $\alpha/(m-j+1)\le\alpha/m_0$. A true
hypothesis is rejected only if the one at rank $j$ is, which requires $p_{(j)}\le\alpha/(m-j+1)\le\alpha/m_0$, so
some $p_i$ with $i\in I_0$ is at most $\alpha/m_0$. By the union bound this has probability at most
$m_0\cdot\alpha/m_0=\alpha$.
\end{proof}

At 10\% one predictor passes on its own, the Treasury bill\ix{Treasury bill} rate with $p=\ovPmin$. Holm's first threshold is
$0.10/14=\ovHolmFirst$, so Holm rejects \ovHolmTen, Bonferroni \ovBonfTen{} and the rule of
\textcite{benjamini1995} \ovBhTen. At 5\% no test passes even alone. Corrections for multiple tests cannot
create evidence; here the single-predictor results do not meet the corrected thresholds, which does
not show that no predictor has genuine forecasting power.
\textcite{harvey2016} argue for thresholds of this kind on new factors in the cross-section.

\subsection{Testing the best of many rules}
\label{sec12:snoop}

Holm\ix{Family-wise error rate}'s procedure and the rule of \textcite{benjamini1995} take the p-values\ix{p value} as given. When the tests are
of trading rules\ix{Trading rule} run over one shared history, the p-values are strongly dependent\ix{Dependence}, and a test built on
their joint distribution can keep its error rate while rejecting more. Let $d_{k,t}$ be the return of
rule $k=1,\dots,m$ in month $t=1,\dots,n$ over a benchmark\ix{Benchmark}, here cash, with mean $\mu_k$ and sample mean\ix{Sample mean}
$\bar d_k$. The question is whether any rule beats the benchmark, so the overall null hypothesis\ix{Hypothesis test} is
$H_0:\max_k\mu_k\le0$, and the hypothesis for rule $k$ is $H_k:\mu_k\le0$.

Each procedure below resamples the $n\times m$ matrix of the $d_{k,t}$ by whole months, which keeps the
dependence across rules, and in blocks\ix{Block bootstrap}, which keeps some of the dependence over time. The stationary\ix{Stationarity}
bootstrap\ix{Bootstrap}\ix{Stationary bootstrap}\index{Stationary bootstrap} of \textcite{politis1994} draws blocks that start at random months and have geometric lengths with mean $b$,
wrapping round the end of the sample. With $B$ resamples it gives bootstrap means $\bar d^{*}_{k,j}$,
$j=1,\dots,B$, and the bootstrap standard deviation $\hat\omega_k$ of $\sqrt n\,\bar d_k$.

\paragraph{The reality check}
\textcite{white2000} tests $H_0$ with $V=\max_k\sqrt n\,\bar d_k$\index{Reality check}. Its bootstrap values are
$V^{*}_j=\max_k\sqrt n\,(\bar d^{*}_{k,j}-\bar d_k)$; centring each rule at its own sample mean imposes the
case $\mu_k=0$ for every rule, the least favourable one for the test. The p-value is the share of the
$V^{*}_j$ at or above $V$. \textcite{sullivan1999} applied the test to thousands of technical rules on a century of
US stock prices.

\paragraph{Superior predictive ability}
\textcite{hansen2005} changes two things\index{Superior predictive ability}. He studentizes, with $t_k=\sqrt n\,\bar d_k/\hat\omega_k$ and
$T=\max(\max_kt_k,0)$, so that a volatile rule does not dominate the maximum, and he centres the bootstrap
at $\hat\mu_k=\bar d_k\mathbf 1\{t_k\ge-\sqrt{2\ln\ln n}\}$, so that a rule that is clearly worse than the
benchmark no longer raises the critical value\ix{Critical value}. The bootstrap statistic\ix{Test statistic} is
$T^{*}_j=\max\bigl(\max_k\sqrt n\,(\bar d^{*}_{k,j}-\hat\mu_k)/\hat\omega_k,\,0\bigr)$, and its p-value is the share of
the $T^{*}_j$ at or above $T$. Centring at $\bar d_k$ instead gives his upper p-value, and centring at
$\max(\bar d_k,0)$ his lower one.

\paragraph{Stepwise testing}
The reality check\ix{Reality check} and the SPA test\ix{Superior predictive ability} ask whether the best rule beats the benchmark. The StepM\ix{Stepwise multiple testing}\index{Stepwise multiple testing} of
\textcite{romano2005} asks which rules do, and controls the familywise error. With
$t^{*}_{k,j}=\sqrt n\,(\bar d^{*}_{k,j}-\bar d_k)/\hat\omega_k$, its first step rejects every $H_k$ whose $t_k$
exceeds $c_1$, the $1-\alpha$ quantile\ix{Quantile} over $j$ of $\max_kt^{*}_{k,j}$. Each later step recomputes the
quantile over the rules not yet rejected, rejects those above it, and the procedure stops at the first
step that rejects nothing.

\paragraph{False discoveries under any dependence}
The rule of \textcite{benjamini1995} controls the false discovery rate\ix{False discovery rate} for independent tests and for some
forms of positive dependence. \textcite{benjamini2001} show that running the same rule at level $q/c(m)$,
with $c(m)=\sum_{i=1}^m1/i$, controls the false discovery rate at $q$ whatever the dependence between the
tests. The price is power: for $m=\mtRules$, $c(m)=\mtCm$.

\begin{proposition}[Ordering of the bootstrap tests]\label{prop12:snoop}
Computed from the same bootstrap draws, Hansen's lower, consistent and upper p-values satisfy
$p_l\le p_c\le p_u$. The critical values of the StepM do not increase from one step to the next, so a
hypothesis rejected at one step stays rejected and the procedure ends after at most $m+1$ steps.
\end{proposition}

\begin{proof}
For each rule the three centring values are ordered, $\max(\bar d_k,0)\ge\hat\mu_k\ge\bar d_k$: if
$\bar d_k\ge0$ the first two are equal; if $\bar d_k<0$ and $t_k$ is at or above the threshold,
$\hat\mu_k=\bar d_k$; and below it $\hat\mu_k=0>\bar d_k$. Each $T^{*}_j$ decreases as the centring
values increase, while $T$ is the same for all three, so the share of draws with $T^{*}_j\ge T$ is
smallest for the lower centring and largest for the upper one. In the StepM the set of rules not yet
rejected only shrinks, so for each draw the maximum of $t^{*}_{k,j}$ over that set can only fall, and so
can its quantile. A statistic above one critical value is above every later one. Each step that does not
stop rejects at least one new hypothesis, which bounds the number of steps.
\end{proof}

\paragraph{Trend rules in managed futures}
Take the six markets of the trend programme of Chapter~9 whose data the repository ships: WTI\ix{Crude oil}\ix{WTI crude oil} crude oil,
the US and Canadian 10 year government bonds\ix{Sovereign bond}, and the euro, the yen and the Canadian dollar against the
US dollar. Returns are monthly, from month-end prices, in excess of the bill rate, from \mtFirst{} to
\mtLast, \mtN{} months. The oil and currency returns are spot returns, without the roll yield or the
interest differential that a futures position would earn, so they differ from what a programme would
have made. A rule is long a market when its cumulative excess return is above its moving average over
$L$ months and short when it is below, for $L$ from \mtLookMin{} to \mtLookMax, which gives \mtRules{}
rules. A change of position costs \mtCost\% of the position each way; the average rule changes position
\mtTurn{} times a year. The bootstrap has $B=\mtB$ resamples with a mean block length of \mtBlock{} months.

Of the \mtRules{} rules, \mtPos{} have a positive mean. The highest mean belongs to the rule on
\mtBestAsset{} with $L=\mtBestLook$, \mtBestMean\% a year, and the highest studentized statistic,
\mtTopT, to the rule on the \mtTopAsset{} with $L=\mtTopLook$, whose bootstrap p-value taken alone is
\mtTopP. One at a time, \mtNaiveFive{} rules pass at 5\% and \mtNaiveTen{} at 10\%. Once the search is
counted, the reality check gives $p=\mtRcP$, the SPA test $p=\mtSpaP$ (lower \mtSpaPl, upper \mtSpaPu),
and the StepM rejects \mtStepM{} hypotheses at 5\%. On the same individual p-values, Holm rejects
\mtHolmFive, the rule of \textcite{benjamini1995} \mtBhFive{} and its version for any dependence
\mtBhyFive; at 10\% the three reject \mtHolmTen, \mtBhTen{} and \mtBhyTen. The first critical value of the
StepM is \mtCrit, against 1.645 for a single one-sided test at 5\% (\cref{fig12:snoop}).

\begin{figure}[t]
  \centering
  \includegraphics[width=0.9\textwidth]{ch12/fig22_snooping}
  \caption{\kind{historical} Studentized mean returns $t_k$ of the \mtRules{} trend rules, \mtFirst{} to \mtLast, by market and
  moving average length. Dotted line: 1.645, the 5\% critical value of one test. Dashed line: \mtCrit,
  the first-step critical value of the StepM at 5\%.}
  \label{fig12:snoop}
\end{figure}

To see what the corrections buy, draw \mtSimReps{} samples of the same length from these excess returns
by the stationary bootstrap, after removing each rule's mean so that no rule beats cash, and run every
procedure on each sample with \mtSimB{} resamples. The best rule passes a 5\% test on its own in
\mtSimNaive\% of the samples. The reality check finds a winner in \mtSimRc\%, the SPA test in
\mtSimSpa\%, the StepM in \mtSimStepM\% and the rule of \textcite{benjamini2001} in \mtSimBhy\%. The last
three are somewhat above the nominal 5\%. They divide by bootstrap estimates of each rule's standard
deviation, which are themselves noisy in fewer than 300 months of heavy-tailed returns; the reality
check does not studentize and stays below 5\% here.

\FloatBarrier
% =====================================================================================
\section{Limitations}
\label{sec12:limits}

The simulations of \cref{sec12:reg,sec12:cv,sec12:trees} use designs chosen to show one effect at a time:
Gaussian predictors, a sparse linear truth or a known nonlinear function, and independent noise. Real
predictors are persistent and the noise heavy tailed, which makes every out-of-sample comparison noisier
than the simulated ones. The network of \cref{sec12:nn} was trained on prices without noise, where the only
error is approximation and optimisation; quoted prices add noise and arbitrage constraints that a network
does not respect unless it is built to.

The equity premium study uses one sample, one evaluation start and one set of tuning rules. Forecasts
starting in another year, other windows for the penalty or other settings for the trees would give other
numbers, and choosing among them would be the search of \cref{sec12:overfit}. The investor's weights are
capped and costs are ignored; with costs, the timing portfolios that trade more would lose more. The series
are those of the update to December 2025, and later updates may revise them.

The hedging study has a known model with constant volatility and no drift, so the learned hedge can be
checked against the delta. Under a misspecified model or with real paths, the regression of
\cref{prop12:qhedge} still learns a hedge, but the hedge is only as good as the simulator, and the intercept is the simulator's
expected payoff under its own measure, a price only when that measure is a pricing measure. The band is optimised over one parameter, and a richer policy
with memory might do better.

\FloatBarrier
% =====================================================================================
\section{Summary}
\label{sec12:summary}

\begin{itemize}
  \item Expected squared error splits into bias, variance and noise. Ridge shrinks each principal direction by
        $d_j^2/(d_j^2+\lambda)$ and, when $X$ has full column rank and the noise variance is positive, always has some small
        penalty whose coefficient error is below that of least squares; the lasso
        soft-thresholds and sets some coefficients to zero. In the sparse example the excess test error falls
        from \bvOlsErr{} for least squares to \bvRidgeErr{} for ridge and \bvLassoErr{} for the lasso
        at the best penalties in hindsight; in a second simulation of the same model, penalties chosen by
        cross-validation leave \cvRidgeErrCv{} and \cvLassoErrCv{} against oracle values of
        \cvRidgeErrOracle{} and \cvLassoErrOracle, keeping most of the gain.
  \item Cross-validation of time series must respect overlap and time: labels that share returns are
        correlated by $(h-k)/h$ and must be purged around each test block. Selecting predictors on the
        full sample and then cross-validating gives an $R^2$ of \kfSelOut\% for a target with no
        predictability.
  \item Bagging and forests average correlated trees, and the variance of the average cannot fall below $\rho s^2$; boosting
        is gradient descent in function space and reduces bias, with the number of rounds and the learning rate as regularisers.
  \item Backpropagation is the chain rule applied layer by layer and costs a small constant multiple of one forward pass. A small
        network learns option prices to a root mean square error of \nnRmse{} basis points of the strike, but its errors and its derivatives need the
        checks of \cref{sec12:nn}.
  \item Out of sample, from \eqOosFrom{} to \eqOosTo, only the Treasury bill rate and the mean of the
        single-predictor forecasts beat the historical mean of the equity premium, with
        $R^2_{\mathrm{OS}}$ of \eqRtwoTbl\% and \eqRtwoComb\%; the forest and boosting did worst of the regularised methods, and least squares on all
        fourteen worst of all, at \eqRtwoOls\%. A small $R^2$ can still matter, by \cref{prop12:ct}.
  \item Quadratic hedging is a regression of the claim on the gains of the features, and recovers the
        Black and Scholes price and delta. With costs, a hedge needs memory: a no-trade band cuts the
        markup by \hdPriceCutWw\%, against \hdPriceCutCost\% for the re-optimised rule without memory.
  \item The best of $N$ trials with no skill has an expected value below $\sqrt{2\ln N}$ standard errors.
        The deflated Sharpe ratio, the probability of backtest overfitting and Holm's procedure all
        discount a result for the search that produced it. On \mtRules{} trend rules in six markets,
        \mtNaiveFive{} pass a single 5\% test, while the reality check ($p=\mtRcP$), the SPA test
        ($p=\mtSpaP$) and the StepM find none that beats cash once the search is counted.
\end{itemize}

\FloatBarrier
% =====================================================================================
\section{Exercises}
\label{sec12:exercises}

\begin{enumerate}[label=\textbf{12.\arabic*}, ref=12.\arabic*, leftmargin=*]
  \item\label{exr12:one} With $X^\top X=nI$ the least squares coefficient of one predictor is 0.8. Find its
  ridge coefficient in \cref{eq12:ridge} with $\lambda=0.25n$ and its lasso coefficient in \cref{eq12:lasso}
  with $\lambda=0.25$. What is the lasso coefficient if the least squares coefficient is 0.2?
  \item\label{exr12:two} An estimate of a mean of 1 is unbiased with variance 0.5. Which multiple $c$ of it has
  the smallest mean squared error, and what is that error?
  \item\label{exr12:three} Trees have variance 1 and pairwise correlation 0.3. Find the variance of a forest
  of 100 trees and of 10.
  \item\label{exr12:four} Labels are twelve month returns of independent monthly returns. Find the
  correlation of labels four months apart. How many labels are unavailable for training around a test block of
  six labels, the block included, with purging and no embargo?
  \item\label{exr12:five} A network has one input $x=1$, one hidden unit $a=\tanh(w_1x)$ with $w_1=0.5$ and an
  output $\hat y=w_2a$ with $w_2=2$; there are no biases. With target $y=0$ and loss
  $\tfrac12(\hat y-y)^2$, find $\hat y$ and the gradients with respect to $w_2$ and $w_1$.
  \item\label{exr12:six} The mean squared error of a model's forecasts is 0.00190 and that of the historical
  mean 0.00192. Find $R^2_{\mathrm{OS}}$.
  \item\label{exr12:seven} The monthly Sharpe ratio of the market is 0.12. By how much does a predictor
  with $R^2=0.5\%$ raise the average excess return of a mean-variance timer?
  \item\label{exr12:eight} In \cref{exr12:six} the mean squared difference between the two forecasts is
  0.00004. Find the mean difference in squared errors and the mean of the Clark and West term $d$.
  \item\label{exr12:nine} A researcher tests 100 independent strategies over five years. Use the
  approximation of \textcite{bailey2014} to find the expected maximum in standard errors and the expected best
  annual Sharpe ratio if none has skill.
  \item\label{exr12:ten} Five p-values are 0.001, 0.008, 0.012, 0.04 and 0.2. At 5\%, how many hypotheses do
  Holm, Bonferroni and the rule of \textcite{benjamini1995} reject?
  \item\label{exr12:eleven} A strategy has a monthly Sharpe ratio of 0.25 over 120 months, skewness $-0.5$
  and excess kurtosis 1. Find the standard error of the ratio and its probabilistic Sharpe ratio against zero.
  \item\label{exr12:twelve} Over one period a claim's change has variance 0.6, the stock's change variance 1.6
  and their covariance 0.9. Find the hedge ratio that minimises the variance of the hedged position and the
  share of the variance it removes.
\end{enumerate}

\subsection*{Computing projects}
\label{sec12:projects}
These use the data and code in the repository. Each is solved by a function in
\texttt{code/ch12/projects.py}, which \texttt{test\_projects.py} runs, so an answer can be
checked with \texttt{python code/ch12/projects.py} from the root of the repository. The
shipped series end in September 2026, and those that end earlier are noted.

\begin{enumerate}[label=\textbf{P12.\arabic*}, leftmargin=*, ref=P12.\arabic*]
\item Rerun the tests of \cref{sec12:snoop} with costs of 0, 10 and 50 basis points per change of position
and mean block lengths of 3, 6 and 12 months. How sensitive are the reality check and SPA p-values to each
choice, and is there any combination in which StepM rejects? Discuss whether such a combination would be
an honest finding.
\item Take the rule with the highest Sharpe ratio among the 66 and compute its probabilistic Sharpe ratio
against zero with \texttt{psr} and its deflated Sharpe ratio for 66 trials with \texttt{dsr}, using the
skewness and excess kurtosis of its returns. Compare the conclusion with that of the bootstrap tests.
\end{enumerate}

\subsection*{Answers}
\begin{enumerate}[label=\textbf{12.\arabic*}, leftmargin=*]
  \item By \cref{prop12:ridge} ridge multiplies by $n/(n+0.25n)$, giving \zaRidge. By \cref{prop12:lasso} the
  lasso gives $S(0.8,0.25)=\zaLasso$, and \zaLassoSmall{} when the least squares coefficient is 0.2.
  \item The error of $c\hat\mu$ is $c^2\times0.5+(1-c)^2$, smallest at $c=1/1.5=\zbC$, where it is \zbMse,
  a third less than the 0.5 of the unbiased estimate.
  \item By \cref{prop12:avg}, $0.3+0.7/100=\zcVar$ and $0.3+0.7/10=\zcVarTen$.
  \item By \cref{prop12:overlap}, $8/12=\zdCorr$. Labels $a-11,\dots,b+11$ are excluded, $b-a+23=\zdPurged$ labels.
  \item $a=\tanh0.5$ and $\hat y=\zeYhat$. Then $\delta_2=\hat y$, $\partial J/\partial w_2=\hat ya=\zeWtwo$ and
  $\partial J/\partial w_1=w_2\hat y(1-a^2)x=\zeWone$.
  \item $1-0.00190/0.00192=\zfRtwo\%$.
  \item By \cref{eq12:ct}, $0.005\times1.0144/(0.0144\times0.995)$, a rise of \zgGain\%.
  \item The model's squared error is lower by \zhGap{} on average, and by \cref{prop12:cw} the mean of $d$
  is that plus the 0.00004, \zhAdj. The adjustment triples the mean difference.
  \item The approximation gives \ziMax{} standard errors, against \ziExact{} exactly, and the best annual
  Sharpe ratio is $\ziMax/\sqrt5=\ziSr$.
  \item Holm's thresholds are 0.01, 0.0125, 0.0167, 0.025 and 0.05, so it rejects \zjHolm; Bonferroni,
  with 0.01 throughout, rejects \zjBonf; the rule of Benjamini and Hochberg, with thresholds 0.01 to 0.05 in
  steps of 0.01, rejects \zjBh.
  \item By Chapter~8 the variance is $(1+0.5\times0.25+\tfrac34\times0.25^2)/119$, so the standard error is
  \zkSe, the statistic \zkZ{} and the probabilistic Sharpe ratio \zkPsr.
  \item $h=0.9/1.6=\zlH$, which removes $0.9^2/(1.6\times0.6)=\zlCut\%$ of the variance, the squared
  correlation.
\end{enumerate}

\FloatBarrier
% =====================================================================================
\section{Notes, sources and credits}
\label{sec12:credits}

\paragraph{Sources}
The bias and variance split, ridge, the lasso, trees, bagging, forests and boosting are covered in
\textcite{hastie2009}. Ridge is from \textcite{hoerl1970}, the lasso from \textcite{tibshirani1996} and its
coordinate descent from \textcite{friedman2010}. Trees are from \textcite{breiman1984}, bagging from
\textcite{breiman1996}, forests from \textcite{breiman2001}, gradient boosting from \textcite{friedman2001},
and the boosting of a smoother from \textcite{buhlmann2003}; the test function of \cref{sec12:trees} is
from \textcite{friedman1991}. Purging and embargoes are from \textcite{lopezdeprado2018}. Backpropagation
is from \textcite{rumelhart1986}, the approximation theorems from \textcite{cybenko1989} and
\textcite{hornik1991}, Adam from \textcite{kingma2015}, and learning option prices with networks from
\textcite{hutchinson1994}. The equity premium study follows \textcite{welch2008} and \textcite{goyal2024},
with the statistics of \textcite{campbell2008} and \textcite{clark2007}, the combination of
\textcite{rapach2010} and the comparison with \textcite{gu2020}. Deep hedging is from
\textcite{buehler2019} and the no-trade band from \textcite{whalley1997}. The probabilistic and deflated
Sharpe ratios are from \textcite{bailey2012,bailey2014}, the probability of backtest overfitting from
\textcite{bailey2017}, and the step-down test from \textcite{holm1979}. The reality check is from
\textcite{white2000} and its use on trading rules from \textcite{sullivan1999}, the test for superior
predictive ability from \textcite{hansen2005}, the StepM from \textcite{romano2005}, the false discovery
rule under dependence from \textcite{benjamini2001}, and the stationary bootstrap from \textcite{politis1994}. The proofs of the propositions are
standard and are given here in the chapter's notation.

\paragraph{Software used}
The methods are written out with NumPy \parencite{numpy} in \texttt{ch12lib.py}, using SciPy
\parencite{scipy} for the normal law, quadrature and L-BFGS (for the hedge weights learned with costs), pandas \parencite{pandas} and openpyxl
\parencite{openpyxl} to read the data, and Matplotlib \parencite{matplotlib} for the figures. The
cross-checks run in R \parencite{rcore} and GNU Octave \parencite{octave}.
\texttt{code/requirements-lock.txt} records the full environment.

\paragraph{Projects for comparison}
The tests compare the tree and the lasso with scikit-learn \parencite{sklearnsw,sklearn2011}, which
is not used to produce any number in the chapter. The full list, with links, is in \texttt{CREDITS.md}.

\paragraph{Data}
The predictors and returns are the monthly data of \textcite{welch2008}, updated by Amit Goyal to December
2025 \parencite{goyalch12}, retrieved on 19 October 2026. They are free to download from his web page but
carry no licence for redistribution, so they are not shipped with the code, and
\texttt{scripts/pull\_ch12.py} downloads them and rebuilds the chapter's data file. The forecasts derived
from them are not shipped either. Every other input is simulated, with the seed in
\texttt{code/ch12/numbers\_ch12.py}.

\end{chapterrefs}
